% !TeX document-id = {f4f6165e-7e15-441d-9a85-fe79757ad5fa}
%======================================================================
%  BACHELOR'S THESIS
%  Ruhr-Universität Bochum — IT-Sicherheit; carried out at MPI-SP
%  Layout: RUB thesis template (bossthesis.cls, bosscoverpage.sty, boss/),
%  adapted for MPI-SP; see the head of bossthesis.cls.
%
%  COMPILE:  latexmk (see latexmkrc): pdflatex -> biber -> pdflatex ...
%            Only the passes an edit actually needs are run.
%            (TeXstudio: F5 / F6 use latexmk via the comments below)
%
% !TeX TXS-program:compile = txs:///latexmk
% !TeX TXS-program:bibliography = txs:///biber
%
%  EDIT config/metadata.tex first. Then write in chapters/.
%======================================================================
\documentclass[twoside,its]{bossthesis}
% options of the class (see bossthesis.cls): [privacy] hides date and place
% of birth; [earlydraft] for quick printouts; [germanthesis] for German.


%----------------------------------------------------------------------
%  Metadata — EDIT THIS FILE, nothing else in config/
%  thesis.tex passes these to the bossthesis class (cover page, title
%  page, declaration). scripts/build_review_docx.py reads the
%  \newcommand lines for the review copy. Open points are inline todos:
%  they sit in the title page and the declaration box, where margin notes
%  cannot go.
%----------------------------------------------------------------------
\newcommand{\thesistitle}{Understanding Privacy Harm Communication in Health Tracking Apps}
\newcommand{\thesiskind}{Bachelor's Thesis in IT Security}
\newcommand{\thesiscitetype}{Bachelor's Thesis}
\newcommand{\authorname}{Hamza Sghaier}
\newcommand{\birthdate}{9. of June 2000}
\newcommand{\birthplacename}{Mahdia, Tunisia}
\newcommand{\advisorname}{Dr.~Kavous Salehzadeh Niksirat}
\newcommand{\examinerone}{Prof.~Dr.~Carmela Troncoso}
\newcommand{\examinertwo}{\todo[inline, inlinewidth=5cm]{Second examiner.}}
\newcommand{\thechair}{Max Planck Institute for Security and Privacy (MPI-SP)}
\newcommand{\thechairaddress}{Max Planck Institute for Security and Privacy\\Universitätsstraße 140 $\cdot$ 44799 Bochum $\cdot$ Germany}
\newcommand{\startdate}{\todo[inline]{Start date set by the examination office.}}
\newcommand{\submissiondate}{\todo[inline]{Submission date.}}

% AI tools named in the declaration box. Drafted from what the repository
% records; the author must check, correct and complete every line.
\newcommand{\aitoollist}{%
\begin{itemize}
 % \item \textit{Claude (Anthropic), used through Claude Code}: writing and
 %   debugging the collection, ranking and checking scripts in
 %   \texttt{scripts/} and the generated tables; preparing the eligibility
 %   screening decisions of the app sample, which the author checks
 %   (Section~\ref{sec:results-sample}); revising and drafting passages of
 %   Chapters~4 and~5 and Appendices~C and~E.
 %   \todo[inline]{Check and complete; name the sections.}
 % \item \textit{OpenAI Codex}: \todo[inline]{The repository contains a Codex
 %   configuration; state whether and for what it was used.}
 % \item \textit{Methods digest of the literature}: \todo[inline]{Name the tool
 %   that produced the cards and summaries in 01\_Literature/AI Literature
 %   Analysis, and its use.}
 % \item \textit{Language model of the coding pipeline}: part of the method,
 %   described in Chapter~\ref{ch:methodology} and Appendix~\ref{app:prompts}.
\item \todo[inline]{AI Usage summary.}
\end{itemize}}



%----------------------------------------------------------------------
%  Packages added to the bossthesis class (RUB template)
%  The class sets fonts, page layout, headings, page styles, captions,
%  hyperref, cleveref, babel and biblatex; this file adds what the thesis
%  needs beyond that and does not load those packages again.
%----------------------------------------------------------------------
\setlength{\emergencystretch}{1.5em}
% Pages may end short instead of stretching the space between paragraphs and
% list items (two-sided report pages are otherwise filled to full height).
\raggedbottom
\hyphenation{Kna-ben-hans}

% --- Colours ----------------------------------------------------------
% Restrained accents for diagrams, source listings and linked references.
% Body text, headings and table rules remain black.
\usepackage{xcolor}
\definecolor{primary}{HTML}{233F5B}
\definecolor{accent}{HTML}{37665D}
\colorlet{accentdark}{accent}
\definecolor{muted}{HTML}{535B63}
\definecolor{linkblue}{HTML}{496D91}

% --- Figures, tables, code ------------------------------------------
% graphicx, booktabs, float, tikz, listings and csquotes come with the class.
\graphicspath{{figures/}}
\usepackage{array}
\usepackage{colortbl}
\arrayrulecolor{black}
\setlength{\heavyrulewidth}{0.8pt}
\setlength{\lightrulewidth}{0.4pt}
\renewcommand{\arraystretch}{1.12}
\usetikzlibrary{arrows.meta,positioning}
% Charts of the literature review (figures/literature-*.tex, generated).
\usepackage{pgfplots}
\pgfplotsset{compat=1.18}
\newcolumntype{L}[1]{>{\raggedright\arraybackslash}p{#1}}
\newcommand{\rowrule}{\cmidrule[0.3pt]{2-4}}
% Figure descriptions (from the proposal's ACM \Description) are not printed;
% they stay in the source for the review copy, which shows TikZ figures as text.
\newcommand{\Description}[1]{}
\usepackage{tabularx}
\usepackage{longtable}
\usepackage{pdflscape} % wide appendix tables on landscape pages
\colorlet{devshade}{primary!12} % rows of a developer with several apps
\captionsetup[table]{position=top, skip=6pt}
\captionsetup[figure]{skip=8pt}
\usepackage{enumitem}
\setlist{itemsep=0.2\baselineskip, topsep=0.4\baselineskip}
\setlist[itemize]{label=\small\textbullet}

% --- Links -----------------------------------------------------------
% The class loads hyperref and cleveref; links keep the thesis's blue.
\hypersetup{
  colorlinks = true,
  allcolors  = linkblue
}

% --- Bibliography ----------------------------------------------------
% The class loads biblatex (Biber, numeric-comp); restrained blue for
% linked citation numbers.
\addbibresource{../_shared/references.bib}
\renewcommand*{\bibfont}{\small}
\setlength{\bibitemsep}{0.35\baselineskip}
% Long URLs: break at letters and digits and set the list ragged right, so
% that lines are not stretched.
\setcounter{biburlnumpenalty}{9000}
\setcounter{biburllcpenalty}{7000}
\setcounter{biburlucpenalty}{8000}
\appto{\bibsetup}{\raggedright}
\DeclareFieldFormat{labelnumberwidth}{\textcolor{linkblue}{\mkbibbrackets{#1}}}

% --- Listings style --------------------------------------------------
% Subtle background and syntax accents; bold keywords also work in greyscale.
\definecolor{codebg}{HTML}{F5F6F7}
\colorlet{coderule}{primary!45}
\colorlet{codekeyword}{primary}
\definecolor{codebuiltin}{HTML}{665078}
\colorlet{codestring}{accent}
\colorlet{codecomment}{muted}
\lstdefinelanguage{json}{
  morestring = [b]",
  morekeywords = {true,false,null},
  sensitive = true
}
\lstdefinestyle{thesiscode}{
  basicstyle      = \linespread{1}\ttfamily\fontsize{8}{10}\selectfont,
  identifierstyle = {}, % the class sets bold blue identifiers
  keywordstyle    = \color{codekeyword}\bfseries,
  keywordstyle    = [2]\color{codebuiltin},
  stringstyle     = \color{codestring},
  commentstyle    = \color{codecomment},
  numberstyle     = \tiny\color{codecomment},
  numbers         = left,
  % Numbers stand outside the left rule with 8pt of clear space before it
  % (absolute units: listings sets em-lengths in the tiny number font).
  numbersep       = 12pt,
  stepnumber      = 1,
  frame           = leftline,
  rulecolor       = \color{coderule},
  backgroundcolor = \color{codebg},
  xleftmargin     = 26pt,
  framexleftmargin = 4pt,
  breaklines      = true,
  breakindent     = 1.2em,
  prebreak        = {},
  postbreak       = \mbox{\textcolor{codecomment}{$\hookrightarrow$}\space},
  columns         = fullflexible,
  keepspaces      = true,
  showstringspaces = false,
  tabsize         = 4,
  aboveskip       = 1ex,
  belowskip       = 1.5ex,
  captionpos      = t
}
\lstdefinestyle{pythoncode}{
  style    = thesiscode,
  language = Python,
  morekeywords = [2]{open,print,len,enumerate,any,dict,str,isinstance,set,list}
}
\lstdefinestyle{jsoncode}{style = thesiscode, language = json}
\lstset{style=thesiscode}

% --- Helpers ---------------------------------------------------------
% Revision marks are the template's todonotes: \todo{...} in the margin,
% \todo[inline]{...} where a margin note cannot go (tables, boxes, the
% title page). Quiet style: pale amber fill, thin amber border and leader
% line, small text, no shadow; the amber matches the orange \highlight
% marks. (size takes a size command only: todonotes also uses it inside
% \csname, so \sffamily there is an error.)
\definecolor{todofill}{HTML}{FFF5E1}
\definecolor{todoline}{HTML}{D9952E}
\definecolor{todotext}{HTML}{4A3410}
\setuptodonotes{color=todofill, bordercolor=todoline, linecolor=todoline,
  textcolor=todotext, size=\scriptsize}
\tikzset{notestyleraw/.append style={line width=0.4pt, rounded corners=1.5pt,
  inner sep=0.7ex}}
\usepackage{xpatch}
\makeatletter
% todonotes redraws the leader line "thick" for every note; make it thin.
\xpatchcmd{\@todo}{thick,}{line width=0.4pt,}{}{%
  \PackageWarning{preamble}{todonotes leader line not restyled}}
\makeatother
\usepackage{soul}

\newif\ifshowhighlights
\showhighlightstrue % Change to \showhighlightsfalse to hide highlights

% A todonotes leader line is drawn from the note's anchor along one text line
% only. A highlighted passage that is followed directly by \todo therefore
% gets a thin bottom edge in the note's colour on every line it spans, which
% runs into the leader line; a highlight without a note stays plain.
% \texthlul draws the highlight and the edge in one soul pass (soul cannot
% nest \hl and \ul); it uses soul's internals (\SOUL@...) as of TeX Live 2026.
\makeatletter
\newdimen\hl@lnht
\newdimen\hl@lnthick \hl@lnthick=0.7pt
\def\hl@ulpreamble{\SOUL@hlpreamble
  \hl@lnht=-\SOUL@uldp \advance\hl@lnht by \hl@lnthick}
\def\hl@ulleaders{\xleaders\hbox to 1pt{%
  \rlap{\vrule width 1.2pt height \SOUL@ulht depth \SOUL@uldp}%
  \rlap{\textcolor{todoline}{\vrule width 1.2pt height \hl@lnht depth \SOUL@uldp}}\hss}}
\DeclareRobustCommand*\texthlul{\SOUL@hlsetup
  \let\SOUL@preamble\hl@ulpreamble \let\SOUL@ulleaders\hl@ulleaders \SOUL@}
\newcommand{\hl@plain}[2]{{\sethlcolor{#1}\hl{#2}}}
\newcommand{\hl@marked}[2]{{\sethlcolor{#1}\texthlul{#2}}}
\newcommand{\hl@off}[2]{#2}
\newcommand{\hl@on}[2]{\@ifnextchar\todo{\hl@marked{#1}{#2}}{\hl@plain{#1}{#2}}}
\newcommand{\highlight}[2][yellow]{%
  \ifshowhighlights
    \expandafter\hl@on % \expandafter closes the conditional before the look-ahead
  \else
    \expandafter\hl@off
  \fi
  {#1}{#2}}
\makeatother



\author{\authorname}
\title{\thesistitle}
\thesistype{\thesiskind}
\thesiscite{\thesiscitetype}
\birthday{\birthdate}
\birthplace{\birthplacename}
\thesisstart{\startdate}
\thesisend{\submissiondate}
\institute{\thechair}
\instituteaddress{\thechairaddress}
\advisors{& \textbf{\advisorname}\\}
\examiners{First examiner: & \textbf{\examinerone}\\
           Second examiner: & \textbf{\examinertwo}\\}
\aitools{\aitoollist}

%======================================================================
\begin{document}

%---------------------------- Front matter ----------------------------
\pagestyle{fancy}
\pagenumbering{roman}

% Cover page, title page and the declaration with the AI-tools box.
\maketitle

\chapter*{Abstract}
\addcontentsline{toc}{chapter}{Abstract}

% Migrated from proposal/proposal.tex (5 October 2026) on 14 October 2026.
Health trackers process data that can reveal sensitive aspects of
a person's life. Research has examined their data practices and users' privacy
concerns, but the way providers describe possible downstream harms needs
closer study. This thesis proposes a content analysis of privacy policies,
store listings and provider websites for 50 widely used apps available in
Germany. Manual walkthroughs will add interface communication.
The analysis will distinguish descriptions of data flows, capabilities such as
inference, and consequences for users. It will examine how harms are presented
and how providers frame privacy protections, guided by literature-based harm
trees.
Coding assisted by a large language model (LLM) will \highlight[orange!15]{be checked against blind
manual coding, with attention to missed mentions. The output will be an
evidence-linked account of harm communication and protection framing in the
captured material, not of data practices, legal compliance or user effects.}\todo{LLM Analysis Methodology is not specified yet}

\noindent\textbf{Keywords:} privacy harms, health trackers, privacy communication, privacy-enhancing technologies, content analysis.



\cleardoublepage
\chapter*{Acknowledgements}
\addcontentsline{toc}{chapter}{Acknowledgements}
\todo[inline]{Do I need this at all?}

\cleardoublepage
\tableofcontents

%---------------------------- Main matter -----------------------------
\cleardoublepage
\pagenumbering{arabic}


% Migrated from proposal/proposal.tex (5 October 2026) on 14 October 2026.
\chapter{Introduction}
\label{ch:introduction}

\section{Problem Statement and Research Gap}
The closest precedent, Malki et al., found inconsistencies between female
health apps' policies and their privacy-related features
\parencite{malki2024exploring}. This thesis instead asks how \emph{health tracking apps} communicate harms across several surfaces, combining harm
coverage, stance and the framing of PETs (Privacy Enhancing Technologies). \highlight[orange!25]{The initial review did not identify any study that combines these elements in this setting. A documented literature search will test and refine this proposed gap.}
\todo{This is a strong claim and should be held until the literature search confirms it with sufficient confidence. The Section should be developped more and more precedent researches to be included.}

\section{Research Questions and Scope}
\noindent\textbf{RQ.} \emph{What types of downstream privacy harms do health
tracker service providers explicitly acknowledge or downplay when describing
their privacy practices or PETs, and how do they frame the protection these
technologies offer?}

\begin{description}[leftmargin=1.5em,itemsep=1pt,topsep=3pt,
                    labelindent=0pt,font=\normalfont\bfseries]
  \item[RQ1] \textit{Coverage.} Which data flows, capabilities and harms are
    mentioned in the captured communication, and which applicable harms
    in the reference model are not found?
  \item[RQ2] \textit{Stance.} How are named harms acknowledged, qualified,
    minimised, denied or presented as the user's responsibility?
  \item[RQ3] \textit{Protection.} What measures and benefits are claimed,
    and how do the scope and strength of these claims relate to the
    reference model?
  \item[RQ4] \textit{Surfaces.} How do coverage, stance and protection claims
    differ between policies, listings, marketing and inspected interfaces?
\end{description}

\noindent\textbf{Scope.} The study covers health tracking apps: apps that let
users record their own health or body data over time, such as activity,
nutrition, sleep, heart rate, mood and reproductive health
(Section~\ref{sec:app-sampling}).
The sample comprises applications from the German storefronts of Google Play and the Apple App Store, as well as the three major applications pre-installed by device manufacturers.
Wearables are considered through their companion apps.

\section{Significance and Expected Contributions}
 \todo{After results/conclusion}

%\noindent\textbf{Users.} A statement about where data is stored may leave the possible
%consequences unclear. Identifying what providers explain can show which
%questions a reader would still need to answer. The US
%post-Roe~\parencite{cao2024deleted} findings motivate attention to fertility data, but do not establish
%the same legal risks for German users.
%
%\smallskip
%\noindent\textbf{Providers and designers.} Comparing surfaces can reveal opportunities
%to explain relevant risks and the limits of protections at decision points.
%
%\smallskip
%\noindent\textbf{Research and policy.}
%\todo[inline]{After Results/Conclusion}

\section{Structure of this Thesis} \todo{After results/conclusion}
%\highlight[orange!25]{Chapter~{\ref{ch:background}} introduces health tracking, privacy harms and
%	the regulatory context. Chapter~{\ref{ch:related-work}} reviews the research
%	that motivates the gap. Chapter~{\ref{ch:methodology}} describes the planned
%	sampling, collection, coding and validation. Chapters~{\ref{ch:results}},
%	{\ref{ch:discussion}} and~{\ref{ch:conclusion}} are reserved for the empirical
%	results, their interpretation and the conclusion. The appendices document
%	the coding scheme, prompts, app sample and worked examples.} 




% Migrated from proposal/proposal.tex (5 October 2026) on 14 October 2026.
\chapter{Background}
\label{ch:background}
\todo{Will be done last.}
%
%In Germany, 73\,\% of smartphone users use health apps to support fitness,
%nutrition or general wellbeing~\parencite{bitkom2025health}. Health trackers can combine information entered by users, such as symptoms
%and cycle dates, with sensor readings and usage records. Together, these data
%may reveal more than users directly provide. A review of 236 studies reports
%that wearable tracker data can support inferences about private attributes,
%including personality traits~\parencite{niksirat2024wearable}.
%
%Data concerning health and a person's sex life are special categories under
%Article~9(1) of the General Data Protection Regulation (GDPR). Activity data
%falls within the health-data definition when
%it reveals information about health status (Article~4(15), Recital~35).
%Articles~5 and~12 require fair and transparent processing and clear information.
%Articles~13 and~14 specify information duties, including purposes, recipients
%and rights~\parencite{gdpr2016}. Guidance allows this information to be layered
%across a policy and interface notices~\parencite{wp292018transparency}. Misleading commercial claims are also subject to
%consumer-protection law~\parencite{ucpd2005}. \highlight[orange!25]{The study does not assume that
%	all communication about harms is voluntary. It asks what providers say about
%	possible consequences alongside their data-practice disclosures.}
%
%A privacy violation and a privacy harm are related but distinct. Citron and
%Solove describe harms such as financial loss, distress and damage to
%relationships~\parencite{citron2021harms}. Knabenhans et al.\ use \emph{harm
%trees} to study the European Digital Identity Framework
%\parencite[\S4.1]{knabenhans2026harms}. A tree links system components, such as
%information disclosed to another party, to capabilities, such as tracking or
%inference, and then to possible harms. Components may arise from the core
%functionality, its digital context or implementation choices. Their analysis
%shows limits of privacy-enhancing technologies (PETs) within that specific
%system and threat model~\parencite{knabenhans2026harms}.
%
%This thesis will adapt the method to trackers. It will distinguish an app's
%intended prediction for its user from an adversary's ability to obtain or use
%that prediction. Martin et al.\ question whether limiting
%third-party access, using inferences instead of raw data, or processing data on
%the device necessarily prevents privacy violations~\parencite{martin2025cookies}.
%The relevant question is which protection is claimed, against which actor and
%under which conditions.



% Migrated from proposal/proposal.tex (5 October 2026) on 14 October 2026.
\chapter{Related Work}
\label{ch:related-work}

\section{Privacy Practices, User Concerns and Risk Communication}
Prior studies examine whether disclosed practices match observed behaviour.
Huckvale et al.\ found that 29 of 36 depression and smoking-cessation apps
transmitted data to Google or Facebook services. Only 12 of 28 apps sending
data to Google and 6 of 12 sending data to Facebook named the respective
company in their policies~\parencite{huckvale2019assessment}. A study of nearly
5{,}000 Android apps found that 81\,\% misrepresented collection or sharing in
the Data Safety section~\parencite{arkalakis2024abandon}. Other work analyses
Apple's privacy labels~\parencite{li2022ios} and fertility trackers through
policy analysis, interface inspection and network measurement
\parencite{shipp2020period,malki2024exploring,hudig2025intimate}.

User studies offer a complementary view. Fitness tracker users considered
harmful data-use scenarios plausible but often unlikely
\parencite{gabriele2020understanding}. US period-tracker users reported privacy
concerns after Roe v.~Wade was overturned~\parencite{cao2024deleted}.
Provider risk communication has also been studied. Zou et al.\ analysed 161
breach notifications drawn at random from those filed with the Maryland
Attorney General between January and June 2018. Hedge
terms such as \emph{may} or \emph{likely} appeared in 112 (70\,\%), and 64
(40\,\%) claimed to have no evidence of misuse, a statement the authors
consider to downplay the risk, because a lack of evidence is not evidence of
absence. Only 9 notifications linked the exposed data to concrete misuse
scenarios, which the authors name as the effective strategy
\parencite[\S5, p.~7]{zou2019might}. Doan et al.\ examined risk and benefit information in
six genetic-testing companies' privacy and research-consent policies
\parencite{doan2024transparency}. Moniz and Almeida gave eight women in
Portugal a design-probe kit on privacy and consent in intimate health
technologies. They conclude that a one-size-fits-all notice leads back to
habituation, and propose that privacy mechanisms should instead inform users
of the specific risks of the data processing, in notices that are tailored to
the context of use, show personalised risk information and use simple language
\parencite[\S5.2]{moniz2025toward}. The authors note that eight participants
may not reflect the target population and that people interested in privacy
are more likely to take part~\parencite[\S3.6]{moniz2025toward}; the proposals
were not tested.

\section{Literature Review}
\todo{To be extended.}



% Migrated from proposal/proposal.tex (5 October 2026) on 14 October 2026.
\chapter{Methodology}
\label{ch:methodology}

%	\section{Draft Status and Study Design}
%	The study plan below is transferred from the proposal dated 5 October 2026.
%	It describes intended work; it does not report a completed study.
%	\todo{Revise this plan against the supervisor meeting of 13 October 2026:
%	remove category quotas; use English primarily with German as fallback;
%	include apps without a privacy policy; extend the keyword space to mental
%	health trackers; exclude connected-device analysis; specify two-coder
%	procedures and justify the sampling and access rules. The ordering and
%	remaining thresholds require a documented decision.}

%%%% Methodology Section Introduction %%%%
This chapter describes how the study will be carried out. The study combines
a structured content analysis~\parencite{krippendorff2019content} with an
adapted cognitive walkthrough~\parencite{wharton1994cognitive}. The sample comprises 50
apps: eligible apps from the German market of Google Play and the App Store,
ranked together by a German-market popularity score based on ratings and reviews, and the health apps that
come preinstalled on the phones of the leading vendor of each platform in Germany. A pilot
with five of these apps will refine the instruments before the main study.
For each app, the privacy policy, the store listings and \highlight[orange!25]{related provider
pages}\todo{Not sure if the providers pages will be collected.} will be collected, and the apps will additionally be inspected in the cognitive walkthroughs.\todo[inline]{Summary of Coding, Analysis, Conclusions, .. after relevant subsections are done} \todo[inline]{Summary of literature review} Figure~\ref{fig:pipeline} summarises this workflow.

\section{Literature Review}
\label{sec:lit-review}
\todo[inline]{Rewritten on 28 October 2026 (Decision D14, which supersedes D13): the search no longer rests on OpenAlex and Semantic Scholar. Check the procedure, and the screening log in 01\_Literature/\allowbreak Literature-Review/, before the supervisor sees it.}
The literature review has three tasks in this study. It collects the harms and
capabilities from which the reference harm trees are built, it supplies the
evidence for the methodological decisions of this chapter, and it tests the
research gap stated in Section~\ref{ch:introduction}. Its procedure follows the
staged review of Baumann et al.~\parencite[App.~A, pp.~56--57]{baumann2025llm}:
fixed Boolean queries built from term groups, the exact string and hit count
reported for every source, works added by hand reported as a source of their
own, removal of duplicates, screening against explicit criteria and the
extraction of fixed fields. Table~\ref{tab:lit-adaptation} lists which of their
choices are kept and which are changed. The search and the screening are reported
with the items of PRISMA~2020 and its search extension
PRISMA-S~\parencite{page2021prisma,rethlefsen2021prisma}. Because the review
maps evidence rather than pooling effects, it is a scoping
review in the sense of Arksey and O'Malley~\parencite{arksey2005scoping}.

\paragraph{Review question.} The review asks one question, derived from the RQ
and the supervisor's statement of the gap, \enquote{what harms service providers
communicate}: \emph{What is known about the privacy harms and capabilities
associated with health tracking apps, about how their privacy communication has
been analysed and what it says, about how language models have been used and
validated for coding privacy text, and about how privacy texts frame risk and
protection?} It has four sub-questions, each tied to a part of this study.
\begin{description}[leftmargin=1.5em,itemsep=1pt,topsep=3pt,labelindent=0pt,font=\normalfont\bfseries]
  \item[LQ1] \textit{Harms and capabilities.} Which privacy harms, risks, threats
    and inference capabilities concern the users of health tracking apps and
    their wearables? It supplies the nodes of the harm trees and the reference
    model of RQ1 and RQ3, and starts from general accounts of privacy
    harm~\parencite{solove2006taxonomy,calo2011boundaries,citron2021harms}.
  \item[LQ2] \textit{Communication.} How has the privacy communication of health
    tracking apps (policies, labels, store listings, notices, settings and
    consent interfaces, marketing) been analysed, and what was found? It informs
    the sampling, material and coding decisions and tests the gap for RQ1--RQ4;
    work on the large-scale analysis of policies is its methodological
    background~\parencite{wilson2016creation,zimmeck2019maps}.
  \item[LQ3] \textit{Language models.} How have large language models been used
    to annotate, classify, extract from or assess privacy text, and how were
    they validated? It informs the LLM-assisted coding.
  \item[LQ4] \textit{Framing.} How do privacy texts in any domain word or present
    risks, harms and protections, for example through vagueness, hedging,
    reassurance or misleading claims? It informs the stance and protection axes
    (RQ2, RQ3).
\end{description}



ZZDIFOPAQUE09820556EF4D686CA30AC2772DAC3B47B945EA6C179101FA3CDCE7A3C9CFA640ZZ



\paragraph{Search concepts.} Eight term groups translate the concepts of the
review question into search terms (Table~\ref{tab:lit-groups}). The domain group
(D) names health tracking apps and their wearables, with the domains of the
sample in Section~\ref{sec:app-sampling} (activity, diet, sleep, mood and
mental health, cycle, fertility and pregnancy) and the generic terms under which
such apps are studied (mHealth, wearable, self-tracking). The privacy group (P)
keeps the domain searches to privacy. The harm group (H) names harms and the
threats and inferences through which capabilities are described. The
communication group (C) names the surfaces the study codes. The privacy-text
group (T) names policies, labels and notices in any domain; the model group (L)
and the task group (A) select language-model studies that annotate or assess
such text; the framing group (F) names the wording strategies that the stance
axis codes. Groups are joined with AND into five strands: S1~=~D AND P AND H
for LQ1, S2~=~D AND P AND C for LQ2, S3~=~T AND L AND A for LQ3, S4~=~T AND F
for LQ4, and S5, the German terms for health apps and data protection, searched
in PubMed only to reach German-language studies of the German market.



ZZDIFOPAQUEBA94E945FAE72D4821F3F03FCB79A8C80E98669371F0ACCE118BC94C49266A22ZZ



\paragraph{Sources.} Venues, where studies are published, are distinguished from
the platforms through which they are searched (Table~\ref{tab:lit-sources-method}).
The venues are those named in the supervisor meeting of 13~October 2026, read
as PoPETs, CCS, USENIX Security, IEEE S\&P and CSCW, together with CHI, where
the supervisor's three example papers appeared, and SOUPS, the venue for usable
privacy; NDSS is added as a reading of the unresolved name \enquote{mdmp}, to be
confirmed. Each venue is listed completely for 2014 to the search date, and the
strands are applied to the listed titles and abstracts. At the security and
privacy venues, privacy is the subject of the venue, so the domain group alone
selects candidates for S1 and S2; at CHI and CSCW the full strands are applied.
Where the platform holds no abstract, titles are matched with two groups
instead of three (for example D AND P), because a title rarely names three
concepts. Four databases add the literature that these venues do not hold:
PubMed for health informatics and medicine, IEEE Xplore for engineering
venues, arXiv for preprints and the ACL Anthology for language processing.
The ACM Digital Library and dblp answered the project's browser and scripts with a
bot verification, which was not completed; Scopus and Web of Science require an
institutional login; Google Scholar offers no export and blocks automated
queries. These sources were therefore not searched, and the ACM venues were
listed through their DOI metadata in Crossref instead. An earlier run of this
review in OpenAlex and Semantic
Scholar~\parencite{priem2022openalex,kinney2023semantic}, on 26 October 2026,
is superseded and not used.



ZZDIFOPAQUE8FB2552A5C2FA882103D2FE62F8D8C7DB060475244C5C28DED4902BBEE6E7340ZZ



\paragraph{Platform strings.} Each strand is translated for each platform, and
the exact strings, fields, filters and search times are given in
Appendix~\ref{app:literature}. PubMed searches titles and abstracts with
truncation inside phrases. IEEE Xplore does not stem, so its strings keep the
truncation and name the title and abstract fields for every term. The arXiv
interface stems words but cannot truncate inside a phrase, and in a test it
returned 34 records for the full S3 string but 62 for \emph{privacy policy} AND
\emph{LLM} alone. Its records are therefore retrieved with short pairs of words
(for example \emph{privacy} AND \emph{fitness}), and the strands are applied to
the retrieved titles and abstracts, as for the venue lists and the ACL
Anthology. No language or publication-type filter is set; both are screening
rules.

\paragraph{Citation searching.} One round of backward and forward citation
searching starts from the four health-tracker papers the supervisor sent on
5~September 2026~\parencite{malki2024exploring,hudig2025intimate,cao2024deleted,niksirat2024wearable}.
Their references and the works citing them are taken from the OpenCitations
Index and their metadata from Crossref.

\paragraph{Existing literature.} Every work cited in the thesis and every paper
of the Methods Digest, the reading collection compiled before the search,
enters the corpus and is screened with the same rules. A thesis citation that
fails a rule is kept in a \emph{contextual collection} with the rule that
excluded it; it may still be cited, for example for a definition or a
statistic, but it is not counted as review evidence.

\paragraph{Deduplication and provenance.} Records that share a DOI, an arXiv
identifier, or a normalised title with publication years at most one year
apart (three years when one of them is an arXiv preprint, which may precede its
paper) are merged into one record. A fuzzy comparison of all titles then
lists near-identical pairs, which are merged by hand when they are the same
work. Each record keeps every route that found it: database, venue, citation
searching, thesis citation or Methods Digest.

\paragraph{Eligibility and screening.} The criteria and the exclusion rules of
Table~\ref{tab:lit-criteria} were fixed before screening. A record is eligible
if it meets C1--C4. Rules R1--R6 are applied to the title, and to the abstract
where the title does not decide; R7--R8 to the full text of the records that
pass. The rules are applied in their order, and an excluded record receives one
primary reason: the first rule that applies. The results therefore report, for
each rule, how many records it assessed and excluded and how many remained.
Every decision and its reason are recorded in a screening log. Where neither
the platform nor Crossref holds an abstract, the abstract is looked up by DOI
in OpenAlex for screening only; OpenAlex is not a search source here. Two
conventions keep the rules apart: systems for hospitals, health records,
telemedicine or the Internet of Things in general fall under R3, and systems
for wearables, body-area networks or health apps under R5; adoption, attitude
and concern studies fall under R4 unless they name specific harms or actors.
\highlight[orange!25]{The screening decisions were prepared by an AI assistant
from titles, abstracts and full texts.}\todo{Check every decision in 01\_Literature/\allowbreak Literature-Review/\allowbreak screening.csv before this section states who screened.}



ZZDIFOPAQUEBAE635C0A3BCF3E78A8648338991B34B09A70DE0C680B60EDAA5A3045B46D92FZZ



\paragraph{Full texts and appraisal.} Full texts are taken from the project's
collection, from open access (publisher, arXiv, PubMed Central) or from the
university library; a record without one is excluded under R7. The review maps
evidence, so no quality score decides inclusion, as PRISMA-ScR allows for
scoping reviews~\parencite{tricco2018prisma}; the peer-review status of each
study is recorded instead, and preprints are marked as such.

\paragraph{Extraction and synthesis.} For every included study, the
sub-questions it answers, the app domain, the surfaces or material studied,
the parts of the RQ it bears on and the harms and capabilities it names are
coded. An evidence map crosses the app domains with the surfaces and with the
parts of the RQ, and shows where evidence is missing. As the counterpart of
Baumann et al.'s Stage~3, the capabilities named in LQ1 studies are collected
in a table that links each capability to its source, the data it needs and the
actor that could use it; this table is the input to the harm trees.

\paragraph{Limits.} One person screens, so no agreement between screeners can
be reported. The ACM Digital Library, Scopus, Web of Science and Google Scholar
were not searched, and five venues could be filtered on titles only, which
lowers recall there; citation searching and the thesis and digest routes
compensate in part. The strands are in English. Later literature is added
through the change log, with a reason, and does not alter the frozen
instruments without a logged revision.



ZZDIFOPAQUEB95F91B077A87C5DCFE2FA00464482D1A624280596D651E81D0F7F0319635F20ZZ



\section{App Sampling}
\label{sec:app-sampling}
The target sample consists of 50 apps. Previous studies that examined both
app documents and the apps themselves used samples of 20 to 64 apps:
Malki et al.\ installed and inspected 20~\parencite{malki2024exploring},
Huckvale et al.\ analysed 36~\parencite{huckvale2019assessment}, and
Hutton et al.\ evaluated 64~\parencite{hutton2018assessing}.

The sample comprises three groups: \textcolor{accent}{(i)} apps preinstalled
by major device manufacturers; \textcolor{accent}{(ii)} apps listed on
Google Play; and \textcolor{accent}{(iii)} apps listed on the Apple App Store.
Preinstalled apps are selected based on the popularity of each manufacturer's
devices. The remaining places are filled with apps from Google Play and the
Apple App Store using defined collection and ranking rules. No quotas are set for the three source groups \textcolor{accent}{(i)}-\textcolor{accent}{(iii)}. The search frame and eligibility rules were
recorded before collection. The resulting sample composition is reported
in Section~\ref{sec:results-sample}.
\todo[inline, color=green]{done}

\paragraph{Domain.} The study concerns health trackers overall. Its operational
scope and search vocabulary draw on two published app studies. Hutton et
al.~\parencite{hutton2018assessing} derived their search categories from the
data types most often reported by 105 members of a Quantified Self group;
their health categories are cycling, diet, exercise, heart, mood, running,
sleep, step count and weight. Malki et
al.~\parencite{malki2024exploring} supply the reproductive-health scope,
including pregnancy tracking, absent from that list. Together these sources
inform the definition and keyword collection for activity, diet and weight,
heart rate, sleep, mood and reproductive health within one health-tracker
domain. Categories describe apps and their features; they do not define
separate sampling, pilot, validation or analysis groups.

\paragraph{Keywords.} The 15 health keywords of Hutton et al.\ and the 12
terms of Malki et al.\ are searched in the original English and in one German
translation each (Table~\ref{tab:keywords}). A translation names the tracked
data type rather than the word (\emph{eating} becomes \emph{Ernährung}).
Identical spellings and repeated translations are searched once, which gives
49 queries.

% Tables 4.1 and 4.2 share one float on a page of their own ([p]).


ZZDIFOPAQUE7B5C09865DFA42A586A671667C2A51F8EE35FDB1DFD62D9333FA3FDC6171ACBDZZ



\paragraph{Search.} Each store is searched in a single session, logged out, in
its German storefront: the App Store through Apple's Search
API~\parencite{apple2022itunessearch} and the Google Play store through its
search page, with the collectors in Appendix~\ref{app:scripts}. For each query,
the first 30 results are kept in both stores. A Google Play search page shows
a single page of results, and deeper results need an interface that the store's
\texttt{robots.txt} disallows. Using the same depth for the App Store follows
O'Loughlin et al., who capped both stores at the number of results Google Play
shows so that neither store is favoured~\parencite{oloughlin2019reviewing}. The
App Store is also searched to 200 results, which serves only to check whether
a deeper search would change the sample. On Google Play, the German count used
for the ranking appears only on an app's product page, so the product page of
every Google Play candidate is opened.

\paragraph{Eligibility.} A listing is eligible if it meets the five criteria
of Table~\ref{tab:eligibility}. Each criterion is checked on the store listing
in the German storefront: title, developer, genre, price and description. An
app listed in both stores is decided on one screened listing. If both of its
listings are screened, the app is excluded when either listing fails a
criterion; in the collection, this never happened. A privacy policy is not
required: an app
without one stays in the sample, and the absence is reported. In-app purchases
and subscriptions are recorded but do not exclude an app. Screening follows the
rank order, and each decision is logged with its reason.

\paragraph{Ranking by German-market popularity.} Candidates are ordered by a
\emph{German-market popularity score}, $P_i$, calculated from the feedback
counts of app~$i$ in the German storefronts: ratings count in the App Store and
written reviews count on Google Play. Feedback counts are used because they are the
only count-based measure that is available for the German market. Download and
install counts are either reported worldwide only or not published at all.
%whereas rating and review counts are shown for each storefront.
The score is
therefore a proxy for popularity used to select the sample; it does not measure
users or downloads. Count-based selection has precedent: Sunyaev et
al.\ selected frequently rated health apps and reported a Spearman correlation
of 0.89 between Android rating counts and downloads~\parencite{sunyaev2015availability}.
Their finding supports the use of feedback counts as a proxy.

The available German-market feedback count differs between stores
(Table~\ref{tab:storefront-fields}):
\begin{itemize}
  \item A \textbf{Google Play listing} contributes its German written-review
    count, denoted $v_i$.
  \item An \textbf{App Store listing} contributes its German rating count,
    denoted $r_i$.
\end{itemize}
Thus, a listing supplies one of the two inputs. The unit being ranked is the
\emph{app}, the same app could be present in both stores, or just in one of both. The scoring formula changes accordingly.
 however: its listings are matched across stores before scoring.
An app found only in Google Play has an observed $v_i$ contribution; one found
only in the App Store has an observed $r_i$ contribution. A matched app found
in both stores has both.

A rating and a written review are different feedback actions. They cannot be
treated as interchangeable counts without a conversion. Google Play reviews
serve as the reference unit: one contributes one score unit. Each App Store
rating contributes $c$ units, where the same conversion factor $c$ is applied
to all candidates. The three cases are therefore
\[
  P_i =
  \begin{cases}
    v_i, & \text{Google Play listing only},\\
    c\,r_i, & \text{App Store listing only},\\
    v_i+c\,r_i, & \text{matched listings in both stores}.
  \end{cases}
\]
The three cases can then be written compactly as
$P_i=v_i+c\,r_i$. The addition combines the converted contributions of the two
listings of a matched app; it does not require every app to have both inputs
or imply that an unobserved store count is actually zero.

The conversion is calculated from the matched candidates found in both stores,
before eligibility screening. Let $M$ contain those candidates with positive
counts in both stores. First, calculate the median ratio of Play reviews to
App Store ratings. Then adjust that ratio using the German platform shares:

\[
  m = \operatorname{median}_{j\in M}\!\left(\frac{v_j}{r_j}\right),
  \qquad
  c = m\,\frac{s_{\mathrm{iOS}}}{s_{\mathrm{Android}}}.
\]

Here, $s_{\mathrm{iOS}}$ and $s_{\mathrm{Android}}$ are the platforms' shares of
German mobile web page views in the month preceding collection, taken from
StatCounter~\parencite{statcounter2026os}. These shares approximate the
platform distribution of users. The observed ratio $m$ reflects both the
different platform sizes and the different rates at which users leave ratings
or reviews. Multiplication by the inverse platform-share ratio aims to remove
the platform-size component, so that $c$ converts App Store ratings into Play
review-equivalent units. This assumes that the typical user split of matched
apps follows the platform shares and that platform-specific feedback rates
are sufficiently similar across apps. 

All candidates are ranked by descending $P_i$, with ties broken by the stable
\emph{candidate key} used by the ranking script.

\paragraph{Selection.} The unit of selection is the provider: a listing is skipped when its developer
is already in the sample, through either store or the preinstalled group. In a
trial collection on 22 October 2026\footnote{Collected with the same collection scripts as in Section~\ref{sec:app-sampling} (Search).}, one developer held 10 of the 47 most-rated
Google Play apps under a single privacy policy (Table~\ref{tab:developer-pilot});
RQ1--RQ3 concern what providers
communicate, so further apps of one provider add little new communication and
reduce the number of independent providers. The rule also means that each app
counts once, and both of its listings are captured.

Two listings belong to the same provider if one of three conditions holds.
(P1) They belong to the same developer account of a store, which both stores
identify by a developer ID\@. (P2) Their developer names agree after case,
accents, punctuation, legal forms and regional entity words are removed, as
\enquote{WW International, Inc.} on the App Store and \enquote{WW
INTERNATIONAL, INC} on Google Play do.
(P3) The names differ, but the listings document one provider: the legal entity
that the App Store shows as seller carries the other developer name, or
listings of the two developer accounts link to the same registered web domain,
as seller or developer website, privacy policy or a link in the description.
Daylio, for example, is published by
\enquote{Relaxio s.r.o.} on the App Store and by \enquote{Habitics} on Google
Play; the App Store seller website, daylio.net, is the domain of the Play
privacy policy. P3 matches are made by hand and recorded with their evidence
(Table~\ref{tab:developer-aliases}). A similar name is no evidence:
\enquote{Leap Health} and \enquote{Leapjoy Games} are different developers.
\todo[inline,color=green]{done!}

\paragraph{Preinstalled apps.} Health apps that ship with the phone reach
users without a download, so their rating counts understate their reach and
they cannot be ranked like the store apps. They form a separate group, and
its members follow from a rule with three steps.
\emph{Vendors:} on each platform, the vendor with the largest share of the
platform's mobile page views in Germany is included. The shares come from the
vendor statistics of StatCounter~\parencite{statcounter2026vendor}, read at the
time of the collection and recorded with their date. A vendor's share of its
platform is its share of all mobile page views divided by the platform's share
of all mobile page views~\parencite{statcounter2026os}, so that each vendor is
compared within its own platform. The rule has no threshold and treats iOS and
Android alike, as the ranking does: a German user of either platform most often
meets the configuration of its leading vendor, the same logic as taking the
most widely used apps. Page views approximate the installed base, not sales.
Rules that admit more vendors, such as a minimum share of all page views or
vendors that together cover most of a platform, are reported as checks.
\emph{Apps:} for each leading vendor, the preinstalled apps are identified from
its documentation and the store listings. An app qualifies if the vendor itself
provides it, if it is installed on the vendor's current phones without a
download, and if it meets the eligibility criteria E1 to E5, in particular E1
(it records the user's own health or body data over time). A vendor that ships
no such app contributes none, and this is recorded.
\emph{Treatment:} the store listing of a qualifying app, if it has one, is
removed from the store rankings, and its vendor counts as a provider already in
the sample, so that the vendor's other apps are skipped by the developer rule.
The vendors and apps found are reported in Section~\ref{sec:results-sample}.

\paragraph{Sensitivity checks.} The selection is repeated with one rule changed
at a time: without the developer rule; with equal weights for the two
platforms; on the alternative scale; with the worldwide Google Play rating count
in place of the German review count; at depth 20; with the App Store at depth
200; with each keyword left out; and with English or German keywords only. It
is also compared with an earlier selection that took apps alternately from a
German App Store ranking and a worldwide Google Play ranking. For each variant,
the number of selected providers that change is reported.

Neither store publishes German figures for every listing field, and the two
stores split their fields differently. Apple shows one summary rating per
country or region, tied to the storefront where the customer bought the
app~\parencite{apple2026ratingsoverview}. Google Play has shown phone users
ratings from their registered country since November
2021~\parencite{grinsted2021making}. Download counts are reported only to an
app's developer, in App Store Connect~\parencite{apple2026acquisition} and in the
Play Console~\parencite{google2026playstatistics}. These sources do not say
which public field follows which rule, so Table~\ref{tab:storefront-fields}
rests on a field check. On 22 October 2026 it requested the same listings from
the German, US and British storefronts and compared the values from Apple's
search and lookup endpoint~\parencite{apple2022itunessearch}, the App Store
product pages, and the Google Play product pages read with an open-source
parser~\parencite{jomingyu2024googleplayscraper}. On the App Store, every rating
figure describes the requested storefront and no worldwide total is published.
On Google Play, the rating count and the install figures were the same in every
storefront, while the average rating and the one-to-five-star breakdown
changed: the breakdown is the country's distribution scaled to the worldwide
count. Both stores show an app's position in a category chart of the
storefront. The check covered few apps in one session, so it shows how the
fields behave and does not enter the sample. The two stores therefore count
ratings on different bases.



ZZDIFOPAQUEC3A475125342E3DD6C1FCDC45948C86352E13333CEE2CBBE012B2CBD49761B6FZZ



\section{Material and Language} For each app, the core corpus will include
the policy, store descriptions, structured store disclosures, the main product
page, and privacy/security pages linked directly from it or the policy. A
fixed supplementary search will capture provider-hosted privacy blog posts
from the preceding 12 months. Ads encountered during collection will form an
opportunistic supplement. The absence of ads or blog posts will not count as
absence of a harm in the core corpus.

Language is recorded per source. English text will be primary for every
captured source, including interface text, and German will be used when no
English version is available; results describe this English-first corpus. If
both versions exist, both will be captured and their differences logged.
Original wording will be coded. The rule does not apply to the search:
keywords combine German and English terms, and the storefronts remain German.
Each source will be archived with URL, date, language, file hash and app
version where applicable. Shared policies will be coded
once and linked to all relevant apps.

\section{App Walkthroughs} Twelve distinct health-tracking apps will be inspected,
six on Android and six on iOS. Selection will follow the
ranking, alternating Android and iOS slots without reusing an app; the balance
serves coverage, not platform comparison. Apps requiring an unavailable paired
device, payment or identity verification will be skipped with a recorded
reason, and unfilled slots reported. Each session will use a fresh install,
synthetic data and a study account where required.

Tasks will cover finding collected data, inferred data, breach consequences
and incident responses, reaching the policy, stopping optional third-party
sharing, exporting data and deleting the account. At each step the researcher
will ask Wharton et al.'s four questions: would the user pursue the intended
goal, notice the action, connect it to that goal and recognise progress
\parencite{wharton1994cognitive}? The pilot will fix starting points and
time limits. Outcomes will distinguish completion, information not found,
unsupported tasks and access barriers. Screen recordings will preserve the
text encountered. This is an expert inspection, not a test of real users.

\section{Reference Harm Trees} Before main coding, reference harm trees will be
built for health trackers from prior research. Sources include a survey of
tracker inferences, threats and vulnerabilities~\parencite{niksirat2024wearable},
users' harm scenarios~\parencite{gabriele2020understanding}, analyses of policies,
data flows and interfaces
\parencite{shipp2020period,malki2024exploring,hudig2025intimate} and post-Roe
user concerns~\parencite{cao2024deleted}. Sources found in the logged search
will be added with a recorded reason. Each pathway will name the data, potential
adversary, access conditions, capability and possible consequence. Required
conditions will be connected explicitly, following Knabenhans et al.'s
AND/OR approach. A prediction made for the user will not itself establish
third-party access (Appendix~\ref{app:example}). As in Knabenhans et al., capabilities include incapabilities, such as a user
being unable to use a feature without agreeing to data sharing.

The component origins will distinguish core functionality, digital context
and implementation choices within the stated model. For every app and
pathway, an applicability record will note the required feature, its
evidence (listing, policy or walkthrough), the access and context
assumptions, the decision (applicable, not applicable or uncertain) and the
reason. A harm category applies to an app when at least one of its pathways
does. The trees are reference models of possible harms, not complete
descriptions of the apps' actual systems. Uncertain cases will be counted and
reported, not silently dropped, so that sparsely documented providers do not
vanish from the gap analysis.

\section{Harm Categories} The starting categories are Knabenhans et al.'s
physical, economic, psychological, autonomy/\allowbreak deprivation, discrimination,
reputational and societal harms~\parencite[App.~A]{knabenhans2026harms},
drawn from Citron and Solove and the CyberPeace Institute. Following them, a
harm is an impairment of the ability to function as one otherwise would in
one's usual context~\parencite[\S4.1]{knabenhans2026harms}; here, the baseline
is the user's situation without the modelled pathway. A relationship tag will
retain Citron and Solove's relationship harms~\parencite{citron2021harms},
which the list omits. New harms found in the pilot will be documented rather
than forced into a category.


\section{Pilot} The pilot takes place after sampling and after the reference
harm trees have been built. Five apps will be drawn at random from the full
health-tracker sample with a recorded seed, without category quotas. The pilot will
test the boundaries of collection and coding and estimate capture time, coding
units, manual effort and model cost. It will also settle the choices left open
in the previous sections: the starting points and time limits of the
walkthroughs, the stance definitions, the handling of new harms, and the LLM
and its prompts. The instruments will then be frozen. The pilot apps will be
recoded with the final instruments and remain in the sample, but will not
enter the validation set.

\section{Coding Units and Axes} To code a text is to attach labels that
record what it says. A unit will be one paragraph or one coherent block of
interface text. Adjacent text will be available for context, but supporting quotations must
be identified precisely. The codebook will define each label, a fitting
example and a counter-example. Table~\ref{tab:axes} summarises the three axes.
Appendix~\ref{app:example} shows the procedure on a constructed example.



ZZDIFOPAQUE23A1B3B675E0E1F7E4C2BEAFB9DF6203A30728C83553E6F8D1DC300530C23C2FZZ



\subsection{Axis A: Coverage} All explicitly mentioned nodes will be
recorded, including whether each is asserted, presented as possible or denied.
The highest level reached will be a separate summary, so a harm
mention does not erase accompanying data-flow information. Feature descriptions can mention inference without presenting it as a risk.

\subsection{Axis B: Stance} Stance labels will be linked to the
particular harm and supporting span, rather than assigned to the paragraph
as a whole. A statement may acknowledge a harm and qualify its likelihood,
and the coder will distinguish uncertainty about a breach from uncertainty
about its consequences. Minimisation will be judged from wording in the unit
or adjacent text, never from inferred intent; qualification alone does not
establish it. RQ2 thus covers stance toward named harms; dismissive language
without one is coded as reassurance. These distinctions adapt Zou et
al.~\parencite{zou2019might} and will be refined in the pilot.

\subsection{Axis C: Protection} The coder will record the measure
and claimed benefit, with actor, data and conditions where specified. Any
technical measure claimed to limit exposure or misuse counts, PET-labelled or
not, including ordinary security controls; non-technical commitments are kept
separate. A \emph{potential mismatch} requires a pathway within the claim's
scope and an explanation of why the measure misses its necessary conditions;
partial protection is not rejected merely because other risks remain. Vague
scope will be recorded as unclear, and no deployed system judged safe or
unsafe.

\section{LLM-Assisted Coding} Polisis and recent LLM approaches are precedents for
data-practice analysis~\parencite{harkous2018polisis,rodriguez2024llm}. The
harm-focused scheme will use prompts and a model selected on the pilot, then
frozen~\parencite{tornberg2024best}; an open-weight model will be preferred if
it performs close to a hosted one, keeping the run repeatable. Prompts will
work in stages: relevant spans, then stance for detected harms, then
protection details. All retained harm mentions and potential mismatches will
be reviewed manually, so the model acts as a screener whose prompts favour
recall. The author will decide tree applicability and mismatches, logging
evidence and reasoning; uncertain outputs or missing spans enter manual
review rather than becoming negative codes. Each unit will run three times;
agreement is a stability check, not validation, and labels returned once are
reviewed rather than discarded. Raw output is kept apart from adjudicated
labels, so performance is not measured after correction.

\section{Validation and Missed Mentions} Ten apps will be drawn at random from
the health-tracker sample, with a recorded seed, from providers and documents not used in
the pilot, stratified by the language of the coded policy where both occur;
the numbers of independent providers and unique documents will be reported.
The author will code them fully before seeing model output, and they will not
be used to tune prompts. Because prompt wording alone can change
conclusions~\parencite{baumann2025llm}, two paraphrased prompts will also be
run and headline counts compared. Disagreements will be adjudicated with each
label's source hidden; two apps will be recoded after at least two weeks,
before drafting, to estimate intra-coder consistency. A second coder will be
sought for applicability, minimisation and mismatch decisions; failing that,
the supervisor will review a sample, including judged non-mismatches, with
disagreements recorded before discussion. A detection is correct only with
the right label, unit and evidence span, and the staged pipeline is scored end
to end, so first-stage misses count. Reported will be per-code precision
(share of model labels supported), recall (share of manual mentions found)
and model--human Cohen's $\kappa$~\parencite{cohen1960coefficient}, with
counts, prevalence and Gwet's AC1, since rare codes can depress $\kappa$
despite high agreement~\parencite{feinstein1990high,gwet2008computing}.
App-level detection is also reported for the validation set as a whole and,
where counts allow, by source and language. Exploratory 95\,\% intervals will resample providers to
respect shared sources; with ten apps they will be wide.

A separate missed-mention audit, reported apart from validation, will search
the remaining corpus for harm-related vocabulary in both languages and check
every unlabelled hit manually. To expose the vocabulary's blind spots, 100
units outside the hits (or all, if fewer), allocated across languages and
sources in advance, will also be read. Recall will be reported as
unestimable where manual positives are too few. For any code with fewer than
20 validation positives or recall below 0.90, all relevant material will be
checked manually or corpus-wide absence claims withheld; these are review
triggers, not proof of accuracy.

\section{Analysis and Reporting} Results will be descriptive. For RQ1, each harm's denominator
will include only applicable apps with the required core sources captured. A
source the provider does not offer is recorded as such; only sources that
exist but cannot be retrieved reduce the denominator. Inapplicable pathways
and text with no identified mention will have different labels. Counts will be per app and surface, with
shared-provider results also shown after grouping repeated providers.
RQ2 will report harm--stance combinations and their denominators. RQ3 will
summarise measures, scope, claim strength and manually reviewed mismatches.
Examples that support and challenge the expected gap will both be retained.

For RQ4, the principal policy--interface comparison will use the same 12
inspected apps; the 50-app document analysis gives the broader overview, and
other surface pairs report how many apps have both surfaces. At the fixed
decision points, onboarding before consent and the privacy-settings screen,
text shown directly, behind a visible link or not found within the task limit
will be distinguished. Inconsistent claims will be reported apart from mere
differences in detail. Absence will mean \enquote{no mention was identified
in the captured material}, with the instrument's measured limitations stated.

\section{Ethics and Feasibility} No human participants are planned.
Walkthroughs use synthetic data, store terms will be checked before automated
retrieval, and any review requirements confirmed before collection. The
pilot's estimates will fix a minimum feasible study (core documents, the blind
validation set, the walkthroughs) before collection; if time runs short,
supplementary marketing collection is reduced first and other changes are
agreed with the supervisor.

\section{Limitations} The sample favours free apps with high German-market
popularity scores found by the chosen keywords and does not represent all
health-tracking apps or their users in Germany; one app per provider also
leaves out further apps of large providers. Feedback counts accumulate over
different app lifetimes and depend on prompts to rate or review. The common
scale assumes a typical platform distribution and approximately stable
platform-specific feedback rates across apps. StatCounter's page-view shares
are an approximation of the user distribution, and a listing not found in one
store contributes no count from that store. These limits prevent interpreting
the score as a count of users or communication encounters. The corpus is a
dated snapshot, walkthroughs cover selected paths, and reference trees depend
on the chosen literature. Unless a second coder joins, one person provides the
reference labels, so model agreement is not independent human-coder
reliability~\parencite{mcdonald2019reliability}. Missed mentions may remain.
The study cannot establish provider intent, actual harms or effects on users.



\chapter{Results}
\label{ch:results}

% Report findings only. Interpretation belongs in the discussion.

\section{Literature Review}
\label{sec:results-lit}
This section reports the search and screening of Section~\ref{sec:lit-review}
and then what the screened literature says about the review question. All
counts come from the search scripts and the screening log; the corpus is listed
in Appendix~\ref{app:lit-inventory}.
\todo{The screening decisions in 01\_Literature/\allowbreak Literature-Review/\allowbreak screening.csv
were prepared by an AI assistant. Check every decision before this section
states who screened, and obtain the 244 full texts still missing.}

\subsection{Search and Screening}

\paragraph{Identification.} The searches ran on 28 October 2026; the final runs
started between 14:46 and 15:10~UTC (Appendix~\ref{app:lit-strings}).
Table~\ref{tab:lit-sources} gives the hits and records of every source. The four
databases returned 2{,}583 hits over the strands, which were 2{,}421 records
once each database's own repeats were merged; IEEE Xplore contributed most of
them (1{,}240), and its S1 query alone returned 1{,}004 hits. The eight venues
were listed completely: 25{,}196 papers, from 372 at SOUPS to 10{,}718 at CHI.
The venue rules selected 85 of them. The German PubMed strand (S5) returned no
record. Merging the database and venue records removed 111 duplicates across
sources and left 2{,}395 records from the searches. The other routes added 585
records that no search had found: 483 from citation searching alone, 64 from the
Methods Digest alone and 38 from the thesis citations alone or from more than
one of these routes.
In total, 2{,}980 records were screened. For 105 of them no abstract was
available, and 371 were screened with an abstract looked up by DOI.



ZZDIFOPAQUE73AFE5FCA9F2EF6FBC1B9FB2EC019950F1A7117571A1E8564F6F5972A46655C3ZZ



Figure~\ref{fig:lit-sources} shows how much each source contributed beyond the
others. Most records of the three large databases were found by that database
alone:
1{,}154 of IEEE Xplore's 1{,}240, 850 of PubMed's 911 and 173 of arXiv's 253.
The overlap is larger where it matters for screening: of the 47 studies that
form the review evidence, 36 are in the Methods Digest and 17 were also found
on arXiv. Citation searching found the most records that passed screening
(90), followed by IEEE Xplore (84), arXiv (69) and PubMed (61). The venue lists
added few records, because the databases and citation searching had already
found most of their papers: CHI contributed 11 records that passed, PoPETs 10,
USENIX Security, SOUPS and IEEE S\&P together 7, CSCW 1, and CCS and NDSS none.



ZZDIFOPAQUE1ACAEDC7685338F113F10942AEB9164E7079AD9AC172AAB1545A9B7CAE7A747EZZ



\paragraph{Screening.} Table~\ref{tab:lit-screening} and
Figure~\ref{fig:lit-exclusions} give the screening by rule, and
Figure~\ref{fig:lit-flow} the whole flow. Each excluded record counts once,
under the first rule that applied, so each rule's denominator is the number of
records left after the rules before it. R3 excluded the most records (1{,}243
of 2{,}502, 49.7\,\%), for example security systems for hospitals, health
records or the Internet of Things in general. R4 removed 585 of 1{,}259 (46.5\,\%) studies of
health apps or wearables without a privacy question, R5 276 of 674 (40.9\,\%)
privacy-preserving systems for wearables and health apps, and R6 106 of 398
(26.6\,\%) studies of privacy text that neither used language models nor
analysed the wording of risk. R1 removed 423 records (14.2\,\%) that were not
research articles, among them 30 chapters and front or back matter of one
edited volume and eight retracted papers or retraction notices; R2 removed 55
works published before 2014. After R1--R6, 292 records passed to full text.



ZZDIFOPAQUE2AC4372B048AED8633B0ABB12F4A6FACF4F9F555375BBDD9C0D7579B5B17F86FZZ





ZZDIFOPAQUE17CC032ECCED4F3DE8B6BA5BAB9D283FBA1535D9FB7FA084E578FF068AC36A48ZZ





ZZDIFOPAQUE57293B6A48A412DBD8301B9F59C0F614A516819AB85A9CF91003E0ADCDC48DEDZZ



\paragraph{Full texts.} Full texts were assessed for the 48 records whose PDF
the project already holds. One was excluded under R8: it discusses the use of
language models for assessing policies without evaluating
them~\parencite{aydin2024assessing}. The other 47 are the review
evidence: 19 answer LQ1, 21 LQ2, 14 LQ3 and 6 LQ4. The full texts of the other
244 records have not been obtained yet; they are open access in most cases
(arXiv, PubMed Central, USENIX, PoPETs) but were not downloaded for this draft.
They remain in the corpus as \emph{awaiting full text}, with the sub-questions
assigned at title and abstract (154 LQ1, 67 LQ2, 37 LQ3, 21 LQ4), and no
statement below rests on them alone.

\paragraph{Existing literature.} All 43 works cited in the thesis before this
review entered the corpus. Sixteen meet the criteria and are review evidence.
The other 27 are in the contextual collection: 10 are not research articles
(R1; web pages, official statistics, a book, a book chapter and software), 5
predate 2014 (R2; among them Solove's taxonomy and Calo's boundaries of privacy
harm), 7 lie outside the domain (R3; among them Citron and Solove's typology,
the harm model of Knabenhans et al.\ and the work on annotation with language
models), and 5 analyse privacy text without language models or risk wording
(R6; OPP-115, Polisis, MAPS, the Data Safety study of Arkalakis et al.\ and the
analysis of Apple's privacy labels by Li et~al.). Five works cited since join
them: the four reporting and review guidelines of
Section~\ref{sec:lit-review} and the paper excluded at full text. Of the 116 Methods Digest
papers, 36 are review evidence, 19 are contextual (cited in the thesis), 60
were excluded and 1 awaits its full text. Citation searching from the four seed
papers yielded 529 DOIs, of which 4 had no metadata; 81 of its 525 records
passed screening, 61 of them found by no other route.

\subsection{What the Literature Shows}

\paragraph{Years and venues.} Figure~\ref{fig:lit-years} shows that the field
is young: of the 291 records that passed screening and were not excluded at full
text, 146 (50\,\%) were published from 2023 onwards, and although S3 searched
from 2022, the first LQ3 records date from 2023. The records come from a wide
range of venues
(Figure~\ref{fig:lit-venues}): 57 from health or medical journals, 51 from
security and privacy venues, 28 from HCI venues and 40 still only as
preprints. Of the 47 studies of the review evidence, 13 appeared in health or
medical journals, 11 at security and privacy venues and 6 at HCI venues; 10
are preprints.



ZZDIFOPAQUE72A07A1A5E1DD3BD1ED3CD06660B6028BB83BF851DF6E049AD2C95EAE22C540FZZ





ZZDIFOPAQUE7CED35FFB76B8E6FF5E084C368FC365B2D236D1136997D126CE597F5ECD09A26ZZ



\paragraph{Evidence map.} Figure~\ref{fig:lit-map-surfaces} crosses the app
domain with the material each study analysed, coded from titles and abstracts
with fixed keyword rules for all 291 records. Wearables in general (90) and
activity and fitness (85) dominate; reproductive health follows with 44,
mental health with 23, and sleep and diet with 9 and 10. Privacy policies are
the most analysed surface (139 records), followed by traffic or code (68) and
interfaces (54). Store listings appear in two records only, one on mental
health apps and one on medical apps, and privacy labels in ten. Of the 88
records answering LQ2, 43 analyse one surface only. The communication of health
tracking apps is thus studied mostly through one document, the policy, and its
relation to app behaviour; the store listing, which every user sees before
installing, is almost unstudied.



ZZDIFOPAQUEBBB805CD240A195B7FBA623B5DF7F5B5C85075757EBF8506C554884F47E68845ZZ



Figure~\ref{fig:lit-map-rq} shows, for the 47 studies of the review evidence,
which part of the thesis RQ each bears on, coded by hand from the full text.
Twenty-six name or test harms (RQ1), ten examine protection claims (RQ3) and
eleven compare several surfaces (RQ4); only five address how harms are worded
or downplayed (RQ2), and of these only the work on privacy
washing~\parencite{brackin2026privacy} also examines protection claims. No study of the review evidence
addresses all four parts, and none codes stance or the framing of PETs in
the communication of health tracking apps. Thirty-five serve as method
precedents.



ZZDIFOPAQUE23B539A5777DD576275A3F9105F937A8CE1040B54AE6AC5E6DC0AC6BDFCB03F5ZZ



\paragraph{LQ1: harms and capabilities.} The review evidence documents
capabilities rather than harms as such: what data an adversary or a recipient
can obtain or infer, and from what. Table~\ref{tab:lit-capabilities} collects
them as the input to the harm trees. Location is inferable from shared activity
routes even when users hide their homes behind the privacy zones that fitness
platforms offer~\parencite{hassan2018analysis,dhondt2022run}; personality traits
are inferable from tracker data~\parencite{zufferey2023watch}; and health and
fitness apps pass data to third parties, often without saying so in their
policies~\parencite{huckvale2015unaddressed,huckvale2019assessment,grundy2019data,tangari2021mobile}.
For reproductive health, studies after the overturn of Roe v.~Wade add law
enforcement as a recipient: users were concerned about access by law
enforcement and third parties~\parencite{cao2024deleted}, and the policies of
female health apps named the possible disclosure of personal data to law
enforcement~\parencite{malki2024exploring}. Mental health apps showed
unnecessary permissions, insecure cryptography and leaks of personal data and
credentials~\parencite{iwaya2023privacy}. The survey of wearable activity
trackers~\parencite{niksirat2024wearable} systematises such risks for
trackers in general. The 154 LQ1 records still awaiting their full text
include many inference attacks from wrist-worn sensors (keystrokes, PINs,
handwriting) and studies of re-identification from activity and heart data;
once assessed, they will extend the table.



ZZDIFOPAQUE075A568841A4EC0B306B65CF211E2784093778E73B735BB3FB4EB9B44D26184CZZ



\paragraph{LQ2: communication.} The 21 LQ2 studies of the review evidence
analyse policies above all: their availability and
quality~\parencite{sunyaev2015availability,robillard2019availability,oloughlin2019reviewing},
their content for menstrual and female health apps~\parencite{shipp2020period,malki2024exploring}
and their match with app behaviour~\parencite{huckvale2019assessment,tangari2021mobile}.
Few go beyond the policy: Hutton et al.\ evaluated apps and policies against
heuristics~\parencite{hutton2018assessing}, Malki et al.\ combined policies
with an inspection of the apps~\parencite{malki2024exploring} and Hudig and
Singh studied the data flows of fertility trackers and users'
preferences~\parencite{hudig2025intimate}. One study rated apps' potential
damage from the information in their store listings; none analyses marketing,
and none codes which harms the communication names. Among the records awaiting
their full text, a discourse analysis of how period- and fertility-tracking
companies describe privacy and an analysis of store description and policy of
mental health apps come closest to this thesis; they must be assessed before
the gap can be stated as confirmed.

\paragraph{LQ3: language models.} The 14 LQ3 studies use language models to
annotate or assess privacy policies, for data practices, contextual
integrity, compliance with privacy law or GDPR transparency, and contradictions
within policies or with labels~\parencite{rodriguez2024llm,chanenson2025automating,xie2025evaluating}.
Each reports an empirical evaluation of the models, as LQ3 required; none of
them annotates harms or stance, which is the coding task of this thesis. Their validation designs are the precedents for the
LLM-assisted coding; the cautions of Baumann et al.\ and Törnberg on
configuration choices and validation remain in the contextual
collection~\parencite{baumann2025llm,tornberg2024best}.

\paragraph{LQ4: framing.} Six studies analyse how texts present risk or
protection: vagueness in policies~\parencite{lebanoff2018automatic}, hedging and
the downplaying of risk in breach notifications~\parencite{zou2019might}, the
transparency of genetic-testing policies~\parencite{doan2024transparency}, the
promises made for privacy-enhancing techniques~\parencite{martin2025cookies},
risk communication in intimate health technologies~\parencite{moniz2025toward}
and privacy washing~\parencite{brackin2026privacy}. Two concern health
(genetic testing and intimate health technologies), and none analyses the texts
of health tracking apps.

\paragraph{Answer to the review question and gaps.} The confirmed evidence
documents capabilities and harms for health tracking apps well, especially
location and attribute inference and third-party sharing, and it shows how
policies of such apps have been sampled and coded. It does not show how
providers word the harms they acknowledge, whether they downplay them, or how
they frame the protection of PETs: no study of the review evidence combines
harm coverage, stance and protection framing across several surfaces of health
tracking apps. The language-model studies validate data-practice annotation,
not the coding of harms. Three limits qualify this answer. First, the 47
studies come disproportionately from the Methods Digest (36), because only
their full texts were at hand; the 244 records awaiting their full text are the
larger and more recent part of the corpus. Second, the domain and surface codes
of the evidence map come from keyword rules on titles and abstracts and are
coarse. Third, the ACM Digital Library, Scopus, Web of Science and Google
Scholar were not searched, and five venues were filtered on titles only.
\section{Sample}
\label{sec:results-sample}
This section reports how the sample of Section~\ref{sec:app-sampling} was
drawn. All figures come from the collector outputs and the ranking script
(Appendix~\ref{app:scripts}); the sample is listed in
Appendix~\ref{app:app-sample}.
\todo{The screening decisions in 04\_Study/\allowbreak App-Sample/\allowbreak screening.csv were
prepared by an AI assistant from the store listings. Check every decision
before this section states who screened.}

\paragraph{Collection.} The App Store was searched on 23 October 2026 between
19:46 and 19:53~UTC\@. Of the 49 queries, 48 returned between 64 and 190 results;
\emph{reproduktive Gesundheit} returned none. The first 30 positions held 1{,}440
hits and 848 distinct listings. Google Play was searched on 24 October 2026
between 15:07 and 15:29~UTC\@. Each query returned between 8 and 50 results, 1{,}178
in total; the first 30 positions held 1{,}138 hits and 645 distinct apps, and
every hit carried an app ID\@. The product pages of all 645 apps were read on
24 and 25 October 2026; one page that timed out was read again, so none is
missing.

\paragraph{Preinstalled apps.} The vendor shares for September 2026, read on
22 October 2026, give Samsung 39.3\,\%, Apple 27.6\,\% and Xiaomi 21.4\,\% of
German mobile page views, and Huawei and Google about 1.5\,\% each
\parencite{statcounter2026vendor}. Apple is the only vendor of iOS (27.7\,\% of
page views). On Android (72.3\,\%), Samsung leads with 54\,\% of the
platform's page views, followed by Xiaomi with 30\,\%
\parencite{statcounter2026os}. The leading vendors are therefore Apple and
Samsung. Their qualifying apps are Apple Health and Apple Fitness on the iPhone
and Samsung Health on Samsung phones. The rule therefore yields three
preinstalled apps, which leaves 47 places for the ranking. A threshold of 10\,\%
of all page views would admit Xiaomi as well, but no preinstalled tracker could
be identified on Xiaomi's international software, so the three apps are the same
under that rule and under a rule that covers 80\,\% of Android. All three have store listings, which
were removed from the ranking; Apple Health, for instance, has 2{,}134 ratings
in the German App Store.

\paragraph{German-market popularity score.} Of the candidates, 125 apps were
found in both stores. The calibration used the 116 with positive rating and
review counts in both stores. Their counts correlated at Spearman 0.79, and
the median ratio was $m = 0.233815$. With the German platform shares of
September 2026, Android 72.3\,\% and iOS
27.7\,\%~\parencite{statcounter2026os}, the conversion factor was
\[
  c = 0.233815\times\frac{0.277}{0.723}\approx 0.089581.
\]
Thus, one App Store rating contributed approximately 0.09 score units and one
Play review contributed one unit. For example, komoot had 417{,}443 German App
Store ratings and 99{,}884 German Play reviews. Using the unrounded factor,
its score was $P_{\mathrm{komoot}}\approx137{,}279$. This is a combined feedback
score, not a user count; the script ranked candidates with unrounded scores.
Across all 644 Google Play candidates, the German review count and the
worldwide rating count correlated at 0.95, yet they rank the leading apps
differently: among the 100 apps with the most installs worldwide, the German
review count ranged from 5 to 99{,}884.

\paragraph{Screening and selection.} Table~\ref{tab:sample-flow} gives the
cascade. Merging the two stores' listings gave 1{,}365 candidates. In descending
order of German-market popularity score, 70 candidates were reached before
the 47 places were filled.
Of these, 11 were excluded at screening because the app records no health or
body data over time (E1), and 12 were skipped because their developer was
already in the sample. Developer names that differ between the two stores were
matched by hand for 13 providers, among them Clue, Pregnancy+, Fddb and
Headspace, each on the evidence listed in Table~\ref{tab:developer-aliases};
no store account was found under two providers, and no two providers among the
70 candidates reached named the same web domain. Three apps whose titles were
translated between the stores were paired by hand. Across the main selection and all sensitivity checks, 275
listings were screened: 156 were eligible and 119 excluded (E1~114, E2~1,
E3~2, E4~1, E5~1).



ZZDIFOPAQUEB463F59D80440369784284156B099BB6C45E754A8805786EF58662C0C34EA7DBZZ



\paragraph{The sample.} The 47 ranked apps come from 47 providers; 33 are listed
in both stores, 8 only on Google Play and 6 only on the App Store. With Apple
Health, Apple Fitness and Samsung Health, the sample holds 50 apps. By primary
category, the ranked apps comprise 22 activity apps, 12 diet and weight apps, 4
cycle trackers, 2 pregnancy apps, 3 sleep apps, 2 mood apps, 1 heart-rate app and
1 wearable companion. The selected apps have between 648 and 447{,}096 German App
Store ratings (39 apps with an App Store listing) and between 1{,}486 and
99{,}884 German Google Play reviews (41 apps with a Google Play listing). All are
free to install, and all 41 Google Play listings name a privacy policy. Of the
47 apps, 42 were found within the first ten positions of at least one query and
5 within positions 11 to 20.

\paragraph{Sensitivity.} Table~\ref{tab:sensitivity} shows how many of the 47
selected providers change when one rule is changed at a time. Without the
developer rule, Leap Fitness Group would take 7 of the 47 places, Simple Design
3, and Google and Nike 2 each (Table~\ref{tab:developer-rule}).



ZZDIFOPAQUE7CF63982303B86265E1D81B6D96AB488EE366694FE5FDE3CEE47BD70E2C24678ZZ



\section{What Privacy Communication Covers}
\todo[inline]{To be written.}

\section{What Privacy Communication Omits}
\todo[inline]{Inference, breaches, incident handling.}

\section{Differences Across Platforms and Categories}
\todo[inline]{To be written.}

\section{Privacy Content in Advertising}
\todo[inline]{To be written.}



\chapter{Discussion}
\label{ch:discussion}

\section{Answering the Research Questions}
\todo[inline]{To be written.}

\section{Implications for Users}
\todo[inline]{To be written.}

\section{Implications for Providers and Designers}
\todo[inline]{To be written.}

\section{Implications for Regulation}
\todo[inline]{To be written.}

\section{Threats to Validity}
\todo[inline]{To be written.}



\chapter{Conclusion}
\label{ch:conclusion}

\section{Summary}
\todo[inline]{To be written.}

\section{Future Work}
\todo[inline]{To be written.}



%---------------------------- Back matter -----------------------------
\cleardoublepage
\pagestyle{fancy-lists}

% Unnumbered lists chapter, as in the template: \phantomsection lets hyperref
% point to the right page, \sectionmark sets the running head.
\phantomsection
\addcontentsline{toc}{chapter}{\listoftitlename}%
\chaptermark{\listoftitlename}

\phantomsection
\addcontentsline{toc}{section}{\listfigurename}%
\sectionmark{\listfigurename}
\listoffigures
\cleardoublepage

\phantomsection
\addcontentsline{toc}{section}{\listtablename}%
\sectionmark{\listtablename}
\listoftables
\cleardoublepage

\phantomsection
\addcontentsline{toc}{section}{\lstlistlistingname}%
\sectionmark{\lstlistlistingname}
\lstlistoflistings
\cleardoublepage

\phantomsection
\addcontentsline{toc}{section}{\bibname}%
\sectionmark{\bibname}
\printbibliography[title={\bibname}]
\cleardoublepage

\pagestyle{fancy}
\appendix

\chapter{Codebook}
\label{app:codebook}
\todo{Full coding scheme with definitions and examples, from codebook v4.}

%----------------------------------------------------------------------
\section{Worked Example}
\label{app:worked-example}

This section walks through the whole procedure once, on one privacy policy,
using codebook version~4. The app, \enquote{CycleMate}, and its policy are
constructed: the text is written so that every important coding decision
occurs at least once. Real policies are expected to stay mostly at the level
of data practices and to name few harms.

\subsection*{How to read this example}

A few terms are used throughout. They are explained here in plain words.

\begin{description}[leftmargin=0pt,itemsep=2pt]
  \item[Unit.] One paragraph of the policy. Coding happens unit by unit.
  \item[Code.] A short label that records what a unit \emph{says}, never what
    the provider might mean. Every code needs a quotation from the unit
    (its \emph{evidence span}).
  \item[Three questions (axes).] \textbf{A}: What is mentioned? \textbf{B}: If
    a harm is named, how is it presented? \textbf{C}: What protection is
    promised?
  \item[Reference tree.] A list of the ways harm \emph{could} happen with a
    health tracker, written down before any policy is read. Each way is a
    \emph{pathway}: data and access conditions, then something an actor
    can do with them (a \emph{capability}), then the harm.
  \item[Applicable.] A pathway is applicable to an app if the app has the
    features the pathway needs. Only applicable harms can be missing from the
    app's communication.
  \item[Potential mismatch.] A protection claim that, read literally, covers
    a pathway which the named measure does not actually block. Only the author
    decides this, never the language model.
\end{description}

\subsection*{Step 0: The reference tree}

The reference trees are built for health trackers from the literature,
following the harm-tree method of Knabenhans et
al.~\parencite{knabenhans2026harms}. \Cref{tab:we-tree} shows five constructed
pathways relevant to CycleMate's features from this shared framework. These
illustrate app-specific applicability rather than a separate study domain. In the
\emph{Conditions} column, AND means that every condition is
needed; OR means that any one of them is enough. Each condition carries its
origin: \emph{core} (the app cannot work without it), \emph{context} (it comes
from running on a phone and network), or \emph{implementation} (the provider
chose it and could choose otherwise).



ZZDIFOPAQUE53E389C000736E75399785CE3E574C59749FAB467C2BC63A88D56F366C7FBD8DZZ



\noindent\textbf{In plain words.} Read each row from left to right: if all
conditions hold, someone gains a capability, and that capability can lead to
the harm. A protection measure blocks a pathway only if it removes one of the
pathway's \emph{necessary} conditions. Three examples:
\begin{itemize}
  \item An app lock removes the access condition of F1, so it can block F1
    even though the cycle data itself is core and stays on the phone.
  \item Turning off backup removes one branch of F2 (the server copy), but
    not the other (a seized phone). F2 stays open.
  \item No measure can remove the logged cycles without stopping the app from
    working. Pathways can only be blocked through their other conditions.
\end{itemize}
The origins are argued from the literature, not proven. The tree describes
\emph{possible} harms; it says nothing about what CycleMate actually does.

\subsection*{Step 1: Which pathways apply to CycleMate?}

Before coding, the author fills in an \emph{applicability record} for each
pathway (\Cref{tab:we-applicability}). It notes the feature the pathway needs,
where that feature is documented, the assumptions made, and a decision:
applicable, not applicable, or uncertain.



ZZDIFOPAQUE4BB7B951BCF04FC94BD7A60441B620B32A3C87C73ACB315C1B376319C301027EZZ



\noindent\textbf{In plain words.} A harm category applies to the app if at
least one of its pathways does. Psychological harm applies (F1, F5),
discrimination applies (F3), reputational harm applies (F5), and so does the
relationship tag (F1). Autonomy harm depends only on F2 and F4, which are both
uncertain, so it is \emph{uncertain}. Uncertain cases are counted and
reported separately; they are never silently dropped.

\subsection*{Step 2: Splitting the policy into units}

The policy is split at paragraph boundaries. A list stays with the sentence
that introduces it. CycleMate's policy gives ten units:

\begin{description}[leftmargin=2.4em,style=sameline,itemsep=2pt]
  \item[U1] \enquote{We collect the information you enter, such as period
    dates, symptoms, sexual activity and mood, to predict your cycle and
    fertile window.}
  \item[U2] \enquote{We share device identifiers and usage data with analytics
    partners to improve our services. We never sell your health data.}
  \item[U3] \enquote{Your health data stays on your device unless you turn on
    backup. This keeps your reproductive health private and protects you from
    anyone using it against you.}
  \item[U4] \enquote{In the unlikely event of a security incident, we will
    notify affected users as required by law.}
  \item[U5] \enquote{If your data were disclosed, it could reveal sensitive
    information such as a pregnancy, which in some jurisdictions could expose
    you to legal risk.}
  \item[U6] \enquote{In the unlikely event that an unauthorised party accessed
    your backup, you might experience unwanted contact.}
  \item[U7] \enquote{Some users worry that a partner could read their entries
    and feel hurt. There is no reason for such concern.}
  \item[U8] \enquote{If you share your cycle calendar with someone, you alone
    are responsible for any embarrassment this causes.}
  \item[U9] \enquote{Because all predictions are calculated with encryption,
    no one can ever use your data to discriminate against you.}
  \item[U10] \enquote{Your data is safe with us. We use industry-standard
    security.}
\end{description}

\subsection*{Step 3: Coding each unit}

\Cref{tab:we-units} gives the codes. In column~A, every mentioned element is
recorded with its status: \emph{asserted} (stated as fact), \emph{possible},
or \emph{denied}. Column~B is filled only when a harm is named, and the stance
is attached to that harm. Column~C records the measure, its scope and how it
is framed.



ZZDIFOPAQUED65B39C6BC7E90711F754525D2189691016008386A33B78E6660FD2C5C41BE05ZZ



\noindent The reasoning behind each row, in plain words:

\begin{description}[leftmargin=0pt,itemsep=3pt]
  \item[U1.] The user enters dates; the app \emph{derives} the fertile window.
    Deriving something the user did not enter is inference (\texttt{K-INF}).
    But the text presents it as a useful feature, not as a risk, and names no
    harm. So there is no stance.
  \item[U2.] The first sentence states a data flow (\texttt{L-CMP}); it never
    says that an analytics partner could build a profile, so \texttt{K-TRK} is
    not coded.
    \enquote{We never sell} is a promise about business practice, not a
    technical measure, so it is coded \texttt{N-COM}, kept separate from
    technical measures.
  \item[U3.] \enquote{Anyone using it against you} is vague: no concrete
    consequence is named, so it is a \emph{generic risk} (\texttt{L-GEN}), not
    a harm, and it gets no stance. The protection claim is on-device storage,
    presented as protecting the user from \enquote{anyone}: harm framing,
    absolute strength. Whether this is a mismatch is decided in Step~4.
  \item[U4.] Promising to notify users \enquote{as required by law} only
    restates a legal duty. It is coded as a general notification commitment
    (\texttt{P-NOT}), not as a concrete incident procedure (\texttt{P-INC}).
    \enquote{Unlikely} creates no stance here, because no harm is named.
  \item[U5.] A concrete consequence is named: legal risk (\texttt{H-AUT}).
    \enquote{In some jurisdictions} narrows \emph{where} the harm applies but
    does not call it unlikely, so the stance is acknowledged (\texttt{S-ACK}).
    Because the harm follows a disclosure, the unit also describes a breach
    consequence (\texttt{P-BRC}).
  \item[U6.] Unwanted contact is named and presented as possible
    (\texttt{S-ACK}). \enquote{In the unlikely event} adds likelihood language,
    so it is also qualified (\texttt{S-QUA}). The qualification targets the
    \emph{breach} (how likely it is), not the \emph{harm} (how bad it would
    be). One harm can carry two stances.
  \item[U7.] The harm (feeling hurt when a partner reads one's entries) is
    named, so it is coded \texttt{H-PSY} with the relationship tag
    (\texttt{T-REL}). It is then dismissed: \enquote{There is no reason for
    such concern} gives no reason in this unit or next to it. That makes it
    minimised (\texttt{S-MIN}). The decision rests on the wording alone, not
    on a guess about the provider's intent.
  \item[U8.] A harm is named (embarrassment), and responsibility for it is
    placed on the user (\enquote{you alone are responsible}): \texttt{S-RSP}.
  \item[U9.] Discrimination is named only to be ruled out (\enquote{no one can
    ever}): its status is \emph{denied}, and the stance is \texttt{S-DEN}. A
    measure is claimed too, so it is also coded under~C.
  \item[U10.] \enquote{Safe with us} and \enquote{industry-standard} name no
    mechanism that could be checked: reassurance without substance
    (\texttt{C-REA}). There is no measure to assess further.
\end{description}

\subsection*{Step 4: The author checks the protection claims}

Each protection claim is compared with the pathways that fall within its
scope (\Cref{tab:we-mismatch}). A \emph{potential mismatch} needs two things:
a pathway that the claim itself covers, and a written reason why the measure
does not remove one of that pathway's necessary conditions. A risk that lies
outside what the claim promises is not enough.



ZZDIFOPAQUE97A7FB19FB04609C8B8C0E65810AF40DA2DFF81F2CFFA5278A6C49194EE1BFCAZZ



\noindent A potential mismatch is a statement about the \emph{text}: what it
promises compared with what the named measure can do. It is not a finding
that CycleMate's actual system is unsafe. The study does not test any app's
security.

\subsection*{Step 5: The result for this policy}

The unit codes are now combined per app and set against the applicability
records from Step~1 (\Cref{tab:we-app}).



ZZDIFOPAQUEA734DE8AD66B4DCCB2B7D7867D786A3256233FDCCEBCD3A64E725206C99C8F0FZZ



\noindent What this one policy contributes to the research questions:
\begin{itemize}
  \item \textbf{RQ1 (coverage).} Four harm categories, including the
    relationship tag, apply to CycleMate. Two are named as possible harms.
    Discrimination is named only in order to be denied. Reputational harm has
    no mention in the captured policy. Autonomy harm is named, but because its
    applicability is uncertain, it is reported separately and neither counted
    as covered nor as missing.
  \item \textbf{RQ2 (stance).} Five named harms carry stances: one
    acknowledged, one acknowledged and qualified, one minimised, one with
    responsibility shifted to the user, and one denied.
  \item \textbf{RQ3 (protection).} Three protection claims are made: a
    no-sale pledge (a practice commitment), on-device storage and encryption.
    The two technical measures are framed as harm prevention with absolute
    strength, and the author flags both as potential mismatches. The no-sale
    pledge is narrowly scoped, and no reference pathway falls within its
    scope. One unit offers reassurance without substance. None of this
    verifies whether the claims are true.
  \item \textbf{RQ4 (surfaces).} The same steps are repeated for the store
    listing, the website and, for the twelve walkthrough apps, the in-app
    screens. The code sets are then compared. For example: does the legal-risk
    warning in U5 also appear on the onboarding screen where the user decides
    whether to sign up?
\end{itemize}
\enquote{No mention identified} always refers to the captured material and to
an instrument with measured limits; it is never a claim that the provider
says nothing anywhere.

\subsection*{Step 6: What the language model does}

The language model helps with Step~3 only, in three stages: it first finds
the relevant spans, then assigns a stance to each named harm, then describes
any protection claim. For every code it must quote the supporting span and
give a short reason; a code without a span is discarded. For U6, the stance
stage returns:



ZZDIFOPAQUEB4EECBDB020E2B0950E1CB1ED1123909ACF24E4027EB6B5E4175B4994E757DE2ZZ



The model never decides applicability (Step~1) or potential mismatches
(Step~4); those are the author's decisions, with the reasoning logged. Every
harm mention the model returns is checked by hand, so the model works as a
screener and is tuned to miss as little as possible. To measure how good it
is, the author codes ten validation apps completely by hand before looking at
the model's output. A detection counts as correct only if the label, the unit
and the evidence span all match. If the first stage misses a span, the codes
that depend on it count as missed as well.



\chapter{LLM Prompts}
\label{app:prompts}
\todo{Every prompt version used, verbatim, with its version number and date.}



\chapter{App Sample}
\label{app:app-sample}
This appendix lists the sample and the evidence for the developer rule
(Section~\ref{sec:app-sampling}). Tables~\ref{tab:app-sample}
to~\ref{tab:developer-rule} follow on landscape pages, and
Table~\ref{tab:developer-aliases} closes the appendix.

\section{Sample}
Table~\ref{tab:app-sample} lists the 50 apps of the sample
(Section~\ref{sec:results-sample}). Numbers 1--47 are the ranked store apps in
descending order of the German-market popularity score; 48--50 are the
preinstalled apps. \emph{Listed in} gives the stores in which the app was found.
$r_i$ is the rating count of the German App Store (23 October 2026, 19:46~UTC);
$v_i$ is the count of written reviews on German Google Play (search of 24
October 2026, 15:07~UTC; product pages of 24 and 25 October); $P_i=v_i+c\,r_i$
with $c=0.0896$ is the popularity score, rounded. A dash means that no listing
from that store is in the search frame.
Categories follow the domain of Section~\ref{sec:app-sampling}. The rows are
generated by \texttt{scripts/app-ranking/v004/rank.py} from the collections in
\texttt{05\_Data/raw/app-sample\_2026-10-25/}.

\section{Developer Rule}
\label{app:developer-rule}
In Tables~\ref{tab:developer-pilot} and~\ref{tab:developer-rule}, shaded rows
belong to a developer that holds more than one of the listed places.

Table~\ref{tab:developer-pilot} lists the 47 Google Play apps with the most
ratings in the trial collection of 22 October 2026 (search and product pages
read between 20:58 and 21:12~UTC; 19 keywords, German storefront, first 20
results; Samsung Health and apps with fewer than 1{,}000 ratings removed). That
collection ranked apps by the worldwide Play rating count and searched Google
Play only, so it has no App Store count $r_i$ and no score $P_i$; $v_i$ is the
German review count of the same pages. The privacy policy is the link on the
listing without its query string. All ten Leap Fitness Group listings link to
\path{leap.app/privacypolicy.html}; their links differ only in a
\texttt{pkg} parameter that names the app. The rows are generated by
\texttt{scripts/sampling-checks/v002/pilot\_table.py}.

Table~\ref{tab:developer-rule} repeats the selection of the sample on the same
collections and screening decisions but without the developer rule: the first
47 eligible candidates in descending order of $P_i$. Its columns are those of
Table~\ref{tab:app-sample}; \emph{Sample no.} gives the app's number in that
table, and a dash marks an app that the developer rule skips because an app of
the same provider ranks higher. The 47 places go to 37 providers. The rows are
generated by \texttt{scripts/app-ranking/v004/rank.py}.

Table~\ref{tab:developer-aliases} lists the developer names that condition P3
matched by hand, with the evidence on the listings. Each row maps the developer
name on the left to the provider on the right.

\begin{landscape}
\begin{singlespace}
\small
\setlength{\tabcolsep}{4pt}


ZZDIFOPAQUED23D0BAAF81060A837DB469228F215C6EB60D417C6204317DC45DBA3D7107CC5ZZ





ZZDIFOPAQUEDF5B35CDFB1D52CD65A297F20996E3F242C04211435275BFC119D3A1AFFC75F0ZZ



\setlength{\tabcolsep}{4pt}


ZZDIFOPAQUE5FA8A143855B694FA48B3B0C2B60704AE49BD62D69D3193197D24F398771873DZZ


\end{singlespace}
\end{landscape}

\begin{singlespace}
\small
\setlength{\tabcolsep}{4pt}


ZZDIFOPAQUEC66F19398D10F66A79CBFB49D869EAD079DA2E96C2717B1DC4E9FD76627BDB96ZZ


\end{singlespace}



% Migrated from proposal/proposal.tex (5 October 2026) on 14 October 2026.
\chapter{Worked Example: From a Pathway to Text Codes}
\label{app:example}

This example applies the method to a constructed privacy notice of a
fictional period tracker, \emph{App~F}; it describes no real app.

\textbf{Step 1: Model the tree.} Figure~\ref{fig:example-path} adapts
Knabenhans et al.'s targeted-advertising tree (their Fig.~6c, with the
tracking sub-tree of Fig.~5b) to App~F~\parencite[\S4.2]{knabenhans2026harms}.
The components keep the paper's roles: C1 is their \emph{reveal other
attributes}, C2 a \emph{persistent identifier}, C3 \emph{online metadata},
which the paper says includes advertising IDs, C4 \emph{repeated
observations}, and C5 the \emph{digitization of data records} that enables
data brokerage. OR means one branch suffices; OR+ means several components
combine. App~F's core function, predicting the cycle, can run on the phone and
reveals nothing to an ad partner, so both branches need implementation
choices that the provider controls and can describe. Dropping the ad SDK
blocks tracking, and not keeping or not selling server records blocks
brokerage; sending interest categories instead of entries does not, because a
category such as \enquote{family planning} still carries the
pregnancy~\parencite{martin2025cookies}.



ZZDIFOPAQUE5EBA084FEBBE6BDC258ABC9065704FB831E4E317E4DB80C5DCAB24D689271CB7ZZ



\textbf{Step 2: Decide whether the tree applies.} The applicability record
notes the required features (C1 and C3 reach ad partners; C5 exists), the
evidence (S2, S3 below), the assumption (the shared events include Pregnancy
Mode, which the text does not state) and the decision: \emph{applicable}, so
psychological harm applies to App~F. Mentions are coded regardless;
applicability only decides which \emph{missing} harms count as gaps (RQ1).

\textbf{Step 3: Code the statements.} Each statement in Table~\ref{tab:worked}
is one unit. A: every mentioned element with its status (asserted, possible,
denied); B: a stance, only for a named harm; C: the measure's framing,
strength and scope. A dash: not applicable.



ZZDIFOPAQUE155961B9C75ADE3BC4C0785ED4BAA2C0CD640C27A56E80951DF6118330E7F6BBZZ



\noindent\textbf{Why these codes.} S1: a prediction made for the user gives
no one else access; no harm is named, so no stance. S2: the data flow is named
but not what partners can do with it, so tracking is not coded. S4: a
capability is denied, not a harm, so no stance. S5: \enquote{rarely}
qualifies how often the harm occurs, so it carries two stances. Dismissing a
capability without naming a harm would be reassurance, not minimisation.

\textbf{Step 4: Review the protection claims.} S3's no-sale pledge covers the
brokerage branch and, taken at its word, closes it: no mismatch. Tracking lies
outside its scope, since C1 and C3 still flow (S2); S3 leaves it open without
contradicting it, which RQ3 records as data-flow framing. S4 promises no
pregnancy-based ads, placing tracking inside its scope; withholding entries
removes neither C1 nor C3, and categories still carry the pregnancy: the
author records a \emph{potential mismatch} with this reasoning. A mismatch
concerns what the text promises, not whether App~F is safe.

\textbf{Step 5: Summarise.} Psychological harm applies and is mentioned once
(S5); tracking is never named, although S2 describes its components. RQ3: one
scoped data-flow claim without mismatch (S3), one absolute harm-framed claim
with a potential mismatch (S4).



% Source of the listings (read from the repository at compile time, not copied):
%   App Store collector   scripts/appstore-scraper/v007/    scraper.py, input.json
%   Google Play collector scripts/googleplay-scraper/v005/  scraper.py, input.json, requirements.txt
% To show a newer version, change the paths below and the two comments above.
\chapter{Scraping Scripts}
\label{app:scripts}
Both collectors read a JSON file and print the result as JSON on standard
output. They record what each store returns for a keyword search. They apply no
eligibility rule and no ranking, and they do not observe app behaviour. The App
Store collector is the version that collected the sample on 23 October 2026,
shown unchanged. The Google Play collector is a shortened form of the version
that collected the sample on 24 and 25 October 2026
(\texttt{googleplay-scraper/v004}, kept in the repository): it drops that
version's page limit, which the sample run did not reach, and gives the same
results from the same search hits and product pages. The collected files are
kept unchanged in \texttt{05\_Data/raw/app-sample\_2026-10-25/}. The inputs
hold the 49 queries of Table~\ref{tab:keywords}.

\section{Input Format}
\begin{description}
  \item[\texttt{keywords}] Search terms, one request each (both scripts).
  \item[\texttt{region}] Two-letter storefront; \texttt{de} selects the German store.
  \item[\texttt{language}] Interface language requested from Google Play.
  \item[\texttt{limit}] App Store: maximum results per keyword (up to 200).
  \item[\texttt{depth}] Google Play: search positions whose apps get their
    product page read (30).
  \item[\texttt{cache}] Google Play: file to which each product page is appended.
\end{description}

\par\medskip\noindent
\begin{minipage}{\linewidth}
\lstinputlisting[style=jsoncode, caption={Input of the App Store collector; the remaining 46 keywords of Table~\ref{tab:keywords} are left out after line~5.}, label={lst:input-apple}, linerange={1-5,51-55}]{../scripts/appstore-scraper/v007/input.json}
\end{minipage}
\par\medskip

\noindent
\begin{minipage}{\linewidth}
\lstinputlisting[style=jsoncode, caption={Input of the Google Play collector; keywords left out as above.}, label={lst:input-play}, linerange={1-5,51-57}]{../scripts/googleplay-scraper/v005/input.json}
\end{minipage}
\par\medskip

\section{App Store Collector}
The script calls Apple's documented iTunes Search API once per keyword and
storefront, keeps iPhone-compatible results in the order returned with their
position before filtering, logs each request with its URL, time, result count and
SHA-256 hash, and pauses between requests to respect the rate limit. Rating counts
and averages describe the requested storefront (Table~\ref{tab:storefront-fields}).

\lstinputlisting[style=pythoncode, caption={App Store collector.}, label={lst:scraper-apple}]{../scripts/appstore-scraper/v007/scraper.py}

\section{Google Play Collector}
The script uses the open-source \texttt{google-play-scraper} parser (version
1.2.7). It first reads the embedded data of the first search page for every
keyword and records each hit's position, title, developer and install band; the
app ID of a highlighted top result is read from its details URL, which the
parser's own field misses. It then opens the product page of every app found
within \texttt{depth}, appending each to the \texttt{cache} file so that an
interrupted run resumes. It restores TLS certificate verification, which the
library disables on import, and rejects any request or redirect that drops the
storefront parameter, so a result cannot silently come from another country.
Failed lookups are recorded under \texttt{errors}. Search positions follow the
order of the search response and are not verified as the store's native
ranking. Rating counts and install figures are worldwide on Google Play; the
one-to-five-star breakdown is the country's distribution scaled to the
worldwide count, and the count of written reviews, which the ranking uses,
follows the storefront country (Table~\ref{tab:storefront-fields}).

\lstinputlisting[style=pythoncode, caption={Google Play collector.}, label={lst:scraper-play}]{../scripts/googleplay-scraper/v005/scraper.py}



% The strings and the table rows are generated by scripts/literature-search/v002/report.py
% (tables/literature-strings.tex, tables/literature-included.tex); change the screening log or the script.
\chapter{Literature Review Records}
\label{app:literature}

\section{Search Strings}
\label{app:lit-strings}
The strings below were sent to each platform on 28 October 2026 (times in
UTC). Bracketed numbers after the arXiv anchor queries are the records each
returned. The strands applied locally to the arXiv pool, the ACL Anthology file
and the venue lists use the term groups of Table~\ref{tab:lit-groups} with the
matching rules of \texttt{lit.py}: text and terms are reduced to lower-case
letters and digits, a trailing * matches any word ending, and terms match only
whole words. The PACM~HCI issues that carry CSCW papers are named
\enquote{CSCW} until 2024; from 2025 the issues are numbered, so all its 2025
and 2026 issues were listed. The search results are kept unchanged in
\texttt{05\_Data/Literature Search/28. Oct/}.

{\small

% Generated by scripts/literature-search/v002/report.py from the frozen search and the screening log; do not edit.
\paragraph{PubMed, S1. Searched 2026-10-28 14:46 UTC; publication date 2014-01-01 to 2026-10-28; 634 records}


ZZDIFOPAQUE1C73BD008159E7D80DB5E63DC7C22A5C56FCACCBE50EB6C15BBF207FDED03B66ZZ



\paragraph{PubMed, S2. Searched 2026-10-28 14:46 UTC; publication date 2014-01-01 to 2026-10-28; 322 records}


ZZDIFOPAQUE321F58B8B2E2BF4093537CA22EA2EC6E4F43454ADDE4F626230EA8AF5B4360EFZZ



\paragraph{PubMed, S3. Searched 2026-10-28 14:46 UTC; publication date 2022-01-01 to 2026-10-28; 13 records}


ZZDIFOPAQUE1713C588BAB000A8A77DA36D2B080CF49973B0DE7487B8175065189BB0CA9EF2ZZ



\paragraph{PubMed, S4. Searched 2026-10-28 14:46 UTC; publication date 2014-01-01 to 2026-10-28; 21 records}


ZZDIFOPAQUEA305808AA3379E6ED0AA678E2049C0CE4065FE3F7D1AE288CA63CB4D40575673ZZ



\paragraph{PubMed, S5. Searched 2026-10-28 14:46 UTC; publication date 2014-01-01 to 2026-10-28; 0 records}


ZZDIFOPAQUEC5D3807507C31F220FDC507726A9A404BCB9AD2E1DA39FB8DD2E98F28498B4D2ZZ



\paragraph{IEEE Xplore, S1. Retrieved 2026-10-28 15:02 UTC; years 2014--2026; 1004 records}


ZZDIFOPAQUE879939A6A465C123B04203C514CEA6917F95FA4A0422F4AB9189613C58944C1EZZ



\paragraph{IEEE Xplore, S2. Retrieved 2026-10-28 15:02 UTC; years 2014--2026; 170 records}


ZZDIFOPAQUE57F53713EFE6CC64150AD21B9BFAB2261B0FC4BAFAE3169E1EA55FB689F69406ZZ



\paragraph{IEEE Xplore, S3. Retrieved 2026-10-28 15:02 UTC; years 2022--2026; 67 records}


ZZDIFOPAQUEB5B868DF2695150CFA3D6A33A8FC1669AC1FD48BD199833BFB11053D6D18B76DZZ



\paragraph{IEEE Xplore, S4. Retrieved 2026-10-28 15:02 UTC; years 2014--2026; 56 records}


ZZDIFOPAQUEF95C1BA02B1CCA76CDF43E03E7A27FE25FA2EBA3EE435BBED870619F5C358FBAZZ



\paragraph{arXiv, anchor queries. Searched 2026-10-28 15:10 UTC; submitted 2014-01-01 to 2026-10-28; 119 queries; pool of 8316 records}


ZZDIFOPAQUE7CFD379655A80658C30148E2C532E72C93580A76159CC6DA2490FC458DD513B5ZZ



\paragraph{Local strands (arXiv pool, ACL Anthology file of 135780 entries, venue lists)}


ZZDIFOPAQUEF0CA4AAE694EA5D83CBE64077DE7BBA36E423D3CAF9321C01508E18A9ACBF3FEZZ




}

\section{Corpus Inventory}
\label{app:lit-inventory}
Table~\ref{tab:lit-included} lists the review evidence, sorted by sub-question,
year and first author; the contextual collection, that is the works cited in
the thesis that do not meet the review criteria, with the rule that excluded
them; and the records that passed title-and-abstract screening but await their
full text, with the sub-questions assigned at that stage. The route codes are
B~database search, V~venue search, C~citation searching, T~thesis citation
before the review and D~Methods Digest. The screening of every
record, with the decision and its reason, is kept in
\texttt{01\_Literature/Literature-Review/screening.csv}.

{\footnotesize
\setlength{\tabcolsep}{3.5pt}


ZZDIFOPAQUE849A71116B6EACAD34BAFA004D148D8DA73A3E667660B2244927911F3EFCE713ZZ


}



\cleardoublepage
\pagestyle{fancy-lists}

% AI Appendix of the RUB template (bossthesis): unnumbered, after the lettered
% appendices. The declaration box at the front refers to it.
\chapter*{AI Appendix}
\addcontentsline{toc}{chapter}{AI Appendix}
\chaptermark{AI Appendix}

This appendix records the use of AI tools for the conceptual design of the
thesis and for its implementation, beyond the list in the declaration.

\section*{AI Use for the Concept}
\todo[inline]{State which AI tools were used for brainstorming or for the design and
conceptual aspects of the thesis, and to what extent; separate own ideas from
AI input.}

\section*{AI Use for the Implementation}
\todo[inline]{State which AI tools were used for which programs and scripts, and to
which end (for example debugging or code generation). The repository's commit
history records Claude Code as co-author of the scripts in \texttt{scripts/}.}



\end{document}
